\chapter{Talcott Parsons \& Robert K. Merton --- Sociological Thinkers}


\section{Talcott Parsons --- The Social System}

\subsection{The Unit Act \& Voluntaristic Theory of Action}

Talcott Parsons' starting point is the \textbf{individual actor} with goals and some means to pursue those goals. Every actor is free to have whatever goal he or she wants in life. The goals and the means, however, are always chosen \textbf{in a context and constrained by norms and situations} --- the \textit{normative} and the \textit{situational}.

\begin{KeyConcept}
\begin{itemize}
\item Parsons tries to show that the actor \textbf{is free} --- unlike the positivists, who say individuals are not free. He clearly says individuals are free to pick up any means to pursue their goals.
\item But the actor is \textbf{equally constrained} by norms and situations. So we cannot really say he is free: he is free to choose means, but the very act of choosing means is constrained by norms and situations.
\end{itemize}
\end{KeyConcept}

\begin{ExampleBox}
\begin{itemize}
\item Suppose I belong to a BPL (below-poverty-line) family and I dream of becoming an IAS officer. The problem is that since I do not have resources, my family is very poor, and I belong to an area where no one pursues civil services as a goal. So my goal, my means, my situation, and my norms may all \textbf{not allow} me to pursue my goal.
\end{itemize}
\end{ExampleBox}

This larger idea --- actor, goals, means, and constraints --- is what Parsons calls the \textbf{unit act}, which he treats as the \textbf{unit of society}: \textit{an actor pursuing goals under certain means and under certain constraints (situational or normative).}

He calls this theory the \textbf{voluntaristic theory of action}. The word ``voluntaristic'' symbolises that the actor \textbf{voluntarily chooses} means and goals, while constraints prevent him from pursuing those goals freely.

The means an actor chooses and the goal he sets for himself can be studied through the \textbf{pattern variables}.

\begin{CriticalPoint}
\begin{itemize}
\item The concept of \textbf{free will} does not really exist in the world (this is also a critique of psychology). Sociology, when you study it, is more about the \textbf{structures that constrain individuals} --- it is largely a \textbf{macro} study.
\end{itemize}
\end{CriticalPoint}

\begin{QuickRevision}
\begin{itemize}
\item Actor $\rightarrow$ goals + means, always constrained by \textbf{norms} (normative) and \textbf{situations} (situational).
\item \textbf{Unit act} = the unit of society: actor + goals + means + constraints.
\item \textbf{Voluntaristic theory of action} = actor voluntarily chooses means/goals, but constraints limit that freedom.
\item Free will is an illusion; sociology studies \textbf{macro structures} that constrain individuals.
\end{itemize}
\end{QuickRevision}

\sectiondivider

\subsection{Pattern Variables (A \& B)}

The \textbf{pattern variable} is a \textbf{model through which we can study the social action of an individual}. The means an actor chooses and the goal he decides can be studied by pattern variable A or pattern variable B.

\begin{ExampleBox}
\begin{itemize}
\item If a Dalit is engaged only in menial tasks, this is studied by \textbf{pattern variable A}.
\item If, despite being Dalit, he is pursuing civil services, this is studied by \textbf{pattern variable B}.
\end{itemize}
\end{ExampleBox}

\textbf{Pattern Variable A} is built on \textbf{ascription} --- the \textbf{ascribed status}, i.e.\ any status you are born with: caste, gender, nationality, religion. \textbf{Pattern Variable B} is built on \textbf{achievement} --- the \textbf{achieved status}, i.e.\ what you achieve rather than are born with.

\begin{ExampleBox}
\begin{itemize}
\item You are born with a gender (e.g., transgender --- this is ascribed status). You become an IAS officer --- that is \textbf{achieved status}. If, as a transgender person, you do only sexual work or begging, you are pursuing the traditional occupations \textit{ascribed} to transgenders (pattern variable A).
\end{itemize}
\end{ExampleBox}

The second contrast is \textbf{particularism vs universalism}.

\begin{ComparisonBox}
\begin{itemize}
\item \textbf{Universalism} = the treatment meted out in the public sphere --- in a classroom, students are just students; no one is special. Officers in the bureaucracy and MPs in administration are treated uniformly. There should be no favoritism or nepotism in the modern world.
\item \textbf{Particularism} = special, unique treatment at home or among one's peer group (those who share your ascriptive features). Favoritism and nepotism symbolise particularism.
\end{itemize}
\end{ComparisonBox}

\begin{CriticalPoint}
\begin{itemize}
\item If your preferences are biased towards ascriptive features --- preferring male over female (as in Parliament), or Hindu over non-Hindu --- you are looking at \textbf{ascription}, not achievement. It shows we are \textbf{not really modern}; in the guise of democracy we are practising nepotism, which is against the principle of democracy.
\end{itemize}
\end{CriticalPoint}

\textbf{How can there be voluntarism along with constraints?}

\begin{ExampleBox}
\begin{itemize}
\item I have a mobile phone in my hand; voluntarily I picked it up (free will). I want to order a pizza, but I cannot because I am in a classroom --- normative constraint (I am a teacher, and should not indulge in such things while teaching).
\end{itemize}
\end{ExampleBox}

The pattern variables also help us \textbf{evaluate and understand social change}.

\begin{DataFact}
\begin{itemize}
\item In ancient/medieval times, recruitment of administrators and warriors was \textbf{not} through competitive exams. People were appointed as advisors because they were \textbf{Brahmin} (even Akbar had Brahmin advisors); the scribes who wrote for rulers belonged to the \textbf{Kayastha} caste (known for beautiful calligraphic writing). Appointment was on the basis of \textbf{caste or religion}.
\item There was a bureaucracy, but it was \textbf{patrimonial, not rational} --- therefore pattern variable A (traditional society).
\item Pre-independence India had no democracy, no constitution, no equal opportunity --- everything fell under pattern variable A. Rulers were decided by religion/heredity (the next Mughal ruler would also be a Mughal).
\end{itemize}
\end{DataFact}

\begin{QuickRevision}
\begin{itemize}
\item \textbf{Pattern variable} = model to study individual social action.
\item \textbf{A} = ascription (ascribed status), particularism --- traditional.
\item \textbf{B} = achievement (achieved status), universalism --- modern.
\item Universalism vs particularism: public-sphere uniform treatment vs home/special treatment (favoritism, nepotism).
\item Pattern variables also track \textbf{social change} (traditional $\rightarrow$ modern).
\end{itemize}
\end{QuickRevision}

\sectiondivider

\subsection{The Concept of System \& Its Three Criteria}

Before defining the social system, we first understand the word \textbf{system}. Something is called a system only when \textbf{three criteria} are met:

\begin{enumerate}
\item \textbf{A set of interrelated parts.} (Example: the human body --- heart interrelated to brain; democracy --- legislature, executive, judiciary interrelated; the state system.)
\item \textbf{These parts work towards the maintenance of equilibrium.} Equilibrium is a peaceful, harmonious state. In the body, all parts adjust to maintain equilibrium. In ecology/environmental science this process is called \textbf{homeostasis}.
\item \textbf{The system resolves conflict in any part to restore equilibrium.} (Example: if blood oozes from the skin, white blood cells automatically heal the wound.)
\end{enumerate}

\begin{ExampleBox}
\begin{itemize}
\item Labour strikes, student protests, farmer protests --- the state intervenes, issues are resolved (through backdoor channels or negotiation), and protests end.
\item MPs and bureaucrats draft bills and send them to stakeholders for opinion (uploaded on the website) so that there is no scope for conflict.
\end{itemize}
\end{ExampleBox}

When these three conditions are met, it is called a \textbf{system}.

\begin{ComparisonBox}
\textbf{Is family a system?}
\begin{itemize}
\item \textbf{Interrelated parts:} mother, father, children.
\item \textbf{Working towards equilibrium:} the family resolves conflict (e.g., if a child is accused of ragging and faces expulsion, the family intervenes to settle the matter).
\item \textbf{Conflict resolution:} any problem with any child is minimised/resolved by the family.
\item $\rightarrow$ Yes, family is a system.
\end{itemize}
\end{ComparisonBox}

\begin{ComparisonBox}
\textbf{Is economy a system?}
\begin{itemize}
\item \textbf{Interrelated parts:} agriculture, manufacturing, services (primary/secondary/tertiary). Formal sector depends on informal (informal provides cheap labour); within formal there is informal (contractual workers, ad hoc professors), and within informal there is formal.
\item \textbf{Working towards equilibrium:} the equilibrium here is \textbf{productivity} --- both sectors work to maximise productivity.
\item \textbf{Conflict resolution:} sexual harassment cells, grievance-addressal mechanisms, managers, HR.
\item $\rightarrow$ Yes, economy is a system.
\end{itemize}
\end{ComparisonBox}

\begin{QuickRevision}
\begin{itemize}
\item Three criteria for a \textbf{system}: (1) interrelated parts, (2) maintenance of equilibrium (homeostasis), (3) conflict resolution to restore equilibrium.
\item Family and economy both qualify as systems.
\end{itemize}
\end{QuickRevision}

\sectiondivider

\subsection{The Social System \& the AGIL Paradigm}

All the organs we see in society --- economy, family, law, media, religion --- are \textbf{social systems}.

\begin{KeyConcept}
\begin{itemize}
\item \textbf{Institutionalisation} or \textbf{constellation} of a large number of social actions by a \textbf{plurality of actors} leads to a social system.
\item Key words: \textbf{constellation}, \textbf{plurality of actors}, \textbf{social action}.
\end{itemize}
\end{KeyConcept}

\begin{ExampleBox}
\begin{itemize}
\item The family as a social system came into existence because a large number of people performed a social action (just as many people performing social action in Weber's Protestant-ethic spirit led to capitalism).
\item People do not have only \textbf{instrumental} needs; they also have emotional/expressive needs --- love, entertainment, appreciation, attention. Different needs give rise to different systems: economy exists for instrumental needs, politics and family for expressive (and instrumental) needs.
\end{itemize}
\end{ExampleBox}

The human body has needs (food, entertainment, love, instrumental needs); society also has a set of needs. When all needs are met, both stay happy. Parsons calls society's needs \textbf{functional prerequisites}.

\begin{KeyConcept}
\begin{itemize}
\item A prerequisite is the minimum eligibility condition before something happens (e.g., to go to America you must have a visa).
\item Society's functional prerequisites must be met for society to ``stay happy'' --- and these are the four needs summarised as \textbf{A G I L} (the \textbf{AGIL paradigm}).
\end{itemize}
\end{KeyConcept}

The four needs and the organs fulfilling them:

\begin{center}
\rowcolors{2}{white}{Tint}
\begin{tabularx}{\textwidth}{p{3.2cm} p{4.2cm} p{3.2cm} Y}
\rowcolor{AccentDark}
\thead{Need} & \thead{Meaning} & \thead{Fulfilled by (organ)} & \thead{Parsons' term for the organ} \\
\midrule
\textbf{A --- Adaptation} & Engaging with the external environment to survive & Economy & \textbf{Organismic system} \\
\textbf{G --- Goal attainment} & Setting and achieving society's goals & Polity & \textbf{Personality system} \\
\textbf{I --- Integration} & Coordinating/integrating all organs & Law, religion, education (many organs) & \textbf{Societal system} (also written ``social system'') \\
\textbf{L --- Latency} & Tension management & Religion, education, family (also media) & \textbf{Cultural system} / \textbf{fiduciary system} \\
\bottomrule
\end{tabularx}
\end{center}

\begin{CriticalPoint}
\begin{itemize}
\item The whole set is a \textbf{social system}; within it there are multiple action systems/subsystems, all \textbf{interrelated}, each working to maintain the larger system (organismic analogy --- society as body).
\end{itemize}
\end{CriticalPoint}

\begin{KeyConcept}
\begin{itemize}
\item \textbf{Fiduciary} means trust. Religion, education, and family are \textbf{fiduciary systems} --- based on trust (trust in parents, trust in God, trust that study leads to a job).
\item \textbf{Trust is the currency of the fiduciary system} = the \textbf{generalised media of exchange}.
\end{itemize}
\end{KeyConcept}

\begin{QuickRevision}
\begin{itemize}
\item Social system = \textbf{constellation} of social actions by a \textbf{plurality of actors}.
\item \textbf{Functional prerequisites} = AGIL (adaptation, goal attainment, integration, latency).
\item A$\rightarrow$economy (organismic), G$\rightarrow$polity (personality), I$\rightarrow$law/religion/education (societal), L$\rightarrow$religion/education/family (cultural/fiduciary).
\item \textbf{Fiduciary} = trust-based; trust is its currency.
\end{itemize}
\end{QuickRevision}

\sectiondivider

\subsection{The Four Action Systems (AGIL) in Detail}

\paragraph{Adaptation (A)}
\textbf{Adaptation} means engaging with the \textbf{external environment} --- you cannot adapt without engaging. It is a need fulfilled by the \textbf{economy} (organismic system), which looks after the \textbf{production and distribution of goods and services} for the \textbf{survival of the community}.

\begin{ExampleBox}
\begin{itemize}
\item In tribal society, the hunting and gathering expeditions were actually the \textbf{economy} playing its role --- not individuals simply going out to hunt. That is why hunting appears in every tribal society: something over and above them (the economy as a system) is making them do it.
\item In modern society, agriculture, manufacturing and services fulfil our survival needs.
\item The \textbf{Indus Valley / Harappan civilisation} grew in the fertile Indo-Gangetic belt --- with canals, drains, irrigation, a bead-making factory and food-processing industries --- the first \textbf{semi-urban} civilisation of India.
\end{itemize}
\end{ExampleBox}

\paragraph{Goal Attainment (G)}
\textbf{Goal attainment} is also a need --- just as an individual has goals (become IAS), society has goals. It is fulfilled by the \textbf{polity} (personality system) --- the political system that \textbf{sets goals}.

\begin{DataFact}
\begin{itemize}
\item Swachh Bharat by 2019; provide a certain number of jobs by 2022 (National Manufacturing Policy); become a net exporter of steel by 2030; sustainable development (SDG); MDG; happiness; ease of doing business.
\end{itemize}
\end{DataFact}

\begin{ExampleBox}
\begin{itemize}
\item Higher ease of doing business $\rightarrow$ more investors $\rightarrow$ more production $\rightarrow$ more employment $\rightarrow$ more GDP.
\end{itemize}
\end{ExampleBox}

The organs that achieve these goals are the \textbf{state bureaucracy} (which executes policies) and the \textbf{corporate sector}. Bureaucracy works through \textbf{hierarchy, chain of command, and division of labour}.

\paragraph{Integration (I)}
\textbf{Integration} is the need for an organ that integrates all the parts (like the brain commands all organs). But in society the brain's function is played by \textbf{many organs} --- not just law, but \textbf{law + religion + education}.

\begin{KeyConcept}
\begin{itemize}
\item The \textbf{constitution} governs the functioning of the legislature, executive, judiciary and constitutional bodies, and imposes limits on each organ so no organ overreaches --- thus integrating all organs. Some things not written in the constitution are interpreted by the judiciary (basic structure).
\item \textbf{Education and religion} play the same role: religion integrates all Hindus as Hindus, all Muslims as Muslims, irrespective of position.
\end{itemize}
\end{KeyConcept}

\begin{CriticalPoint}
\begin{itemize}
\item Parsons formulated this theory taking \textbf{America} as his model --- America in the 1930s--40s was still homogeneous (largely 90\% Catholic/Jewish; few Sikhs or Hindus then).
\end{itemize}
\end{CriticalPoint}

\paragraph{Latency (L)}
\textbf{Latency} means \textbf{tension management}. Just as individuals get tension and need friends/advice/counselling, society also develops tensions, which Parsons calls \textbf{strains}.

\begin{ExampleBox}
\begin{itemize}
\item India feels strain seeing China (once like us) advance rapidly --- the root cause being China's economic growth despite ``illegitimate economic practices'' that the WTO, IMF and World Bank cannot check (the world is now dependent on Chinese goods --- almost 40\% of the world economy).
\item The 1990s balance-of-payments crisis; the coronavirus striking India.
\end{itemize}
\end{ExampleBox}

Strains are managed by \textbf{religion, education and family}:

\begin{ExampleBox}
\begin{itemize}
\item \textbf{Religion:} during corona, Ramayan and Mahabharat were aired, a ``Corona Mata'' mandir was built, people became more religious.
\item \textbf{Education:} scientists educated themselves on the origin of the coronavirus strain $\rightarrow$ could develop a vaccine.
\item \textbf{Family:} the family diffuses tension (takes care of you after a bad result or illness).
\end{itemize}
\end{ExampleBox}

\begin{KeyConcept}
\begin{itemize}
\item The cultural (fiduciary) system \textbf{motivates individuals and diffuses tension}, and also fulfils needs like love, pleasure, entertainment. Media can be placed here (though media also plays an integrative role).
\end{itemize}
\end{KeyConcept}

\paragraph{Structural functionalism (the larger frame)}
This is what \textbf{structural functionalism} is: each organ of society, interrelated with the others, contributes to the functioning of the system as a whole --- a theory of the \textbf{maintenance of society}.

\begin{CriticalPoint}
\begin{itemize}
\item A \textbf{single social action is not a system} --- a system is made of a plurality of actors performing social action.
\item Parsons \textit{did} discuss conflict (one criterion of a system is organs managing conflict), but he treats conflict as something \textbf{managed within the system}, not a big issue --- unlike the \textbf{conflict paradigm}, which gives importance only to conflict.
\item Therefore Parsons cannot explain \textbf{poverty, gender inequality, domestic violence} --- if all organs worked smoothly, why would these social issues exist?
\item Adaptation is related only to the \textbf{economy} (economy is the base, as in Marx). ``Personality sentiment'' is not an organ of society.
\end{itemize}
\end{CriticalPoint}

\begin{QuickRevision}
\begin{itemize}
\item \textbf{Adaptation} = engaging external environment; economy = production + distribution for survival.
\item \textbf{Goal attainment} = polity sets goals; achieved by state bureaucracy + corporate sector.
\item \textbf{Integration} = law (constitution) + religion + education coordinate all organs.
\item \textbf{Latency} = tension management; religion/education/family (fiduciary/cultural system) motivate and diffuse strain.
\item Structural functionalism = theory of society's \textbf{maintenance}; weak on conflict (can't explain poverty, inequality, domestic violence).
\end{itemize}
\end{QuickRevision}

\sectiondivider

\subsection{The Tribal Society: All Organs in Primitive Form}

Even tribal societies had all the AGIL organs, in primitive form:

\begin{center}
\rowcolors{2}{white}{Tint}
\begin{tabularx}{\textwidth}{p{3.5cm} Y}
\rowcolor{AccentDark}
\thead{Organ} & \thead{Primitive (tribal) form} \\
\midrule
Economy & \textbf{Primitive economy} --- hunting, gathering, cultivation \\
Polity & \textbf{Tribal chiefdom} --- the leopard-skin chief (adorned leopard skin as sign of leadership) \\
Religion & \textbf{Animism} (spirit/ancestor worship), \textbf{naturism} (nature worship --- Agni Dev, Vayu Dev, Jal Dev), \textbf{totemism} \\
Education & The art of hunting/gathering --- survival knowledge transferred to the next generation; environmental knowledge (medicinal virtues of plants) \\
Family & A \textbf{community} --- no nuclear/joint division; ``communism of wives''; repressive laws \\
\bottomrule
\end{tabularx}
\end{center}

\begin{DataFact}
\begin{itemize}
\item \textbf{Animism} = worship of spirits/ancestors.
\item \textbf{Naturism} = worship of nature. According to \textbf{Max Müller}, nature worship is the most ancient religion.
\item \textbf{Totemism} --- according to \textbf{Durkheim} (whose tribal society was the \textbf{Arunta} tribe), totemism is the most ancient religion.
\item Collective conscience was sacred; anyone disturbing the moral fabric was executed (repressive law).
\end{itemize}
\end{DataFact}

\begin{KeyConcept}
\begin{itemize}
\item \textbf{Education does not mean school} --- it means the transfer of survival skills (hunting, gathering) and values to the next generation. Tribal education was also very \textbf{environmental} (learning medicinal virtues of plants from childhood).
\item Tribal family had no strict boundary --- ``communism of wives'' (a concept Plato discussed, and which is discussed in Engels' account of the origin of private property, family and state).
\end{itemize}
\end{KeyConcept}

\begin{QuickRevision}
\begin{itemize}
\item Tribal society had all AGIL organs in primitive form: hunting/gathering economy, leopard-skin chiefdom, animism/naturism/totemism, survival-skill education, and community-family.
\item \textbf{Naturism} (Max Müller) vs \textbf{totemism} (Durkheim, Arunta) --- both claimed as ``most ancient religion.''
\end{itemize}
\end{QuickRevision}

\sectiondivider

\subsection{The Cybernetic Hierarchy of Control}

The bigger question is: \textbf{how are these organs interrelated to each other?} To answer this, Parsons gave another theory --- the \textbf{cybernetic hierarchy of control} --- to explain the interrelationship between the \textbf{action systems} (also called subsystems).

\begin{KeyConcept}
\begin{itemize}
\item \textbf{Cybernetics} is basically a \textbf{program of communication} (communication between components). Just as the CPU (the brain of a computer) passes information to monitor/mouse/keyboard, and the brain passes information to all organs, the action systems are interrelated through cybernetic communication.
\end{itemize}
\end{KeyConcept}

The key arrangement is the \textbf{LIGA} order (top to bottom: L $\rightarrow$ I $\rightarrow$ G $\rightarrow$ A):
\begin{itemize}
\item \textbf{L (culture)} and \textbf{I} are \textbf{high in information} (or ``control'').
\item \textbf{G} and \textbf{A (economy)} are \textbf{high in energy}.
\item \textbf{Information flows downward} (culture $\rightarrow$ economy), while \textbf{energy flows upward} (economy $\rightarrow$ culture).
\end{itemize}

\begin{DataFact}
\begin{itemize}
\item Higher systems (\textbf{L, I}) \textbf{control} G and A; and G and A \textbf{pass energy} to L and I.
\item \textbf{Economy is highest in energy} --- like Marx, who believed economy is the base and everything else superstructure (though Parsons does not use the word ``superstructure'').
\item Development in any society largely happens because of \textbf{economic and technological advancement} (industrial revolution $\rightarrow$ advancement of European civilisations).
\end{itemize}
\end{DataFact}

\begin{CriticalPoint}
\textbf{Culture as regulatory authority:}
\begin{itemize}
\item Whatever development happens at the economic/political level \textbf{must be acceptable to the culture}. If culture (the ``brain'') does not accept the change, it will ``vomit'' it.
\item \textbf{Examples:}
\begin{itemize}
\item McDonald's was initially unsuccessful in India, then succeeded with the \textbf{McAloo Tikki} --- selling the same global product by inculcating local practice. \textbf{Microsoft} adopted the Devanagari script/font in Microsoft Word. This ``global + local'' is \textbf{glocalisation}, a term coined by \textbf{Ronald Robertson}. (E.g., Punjabi folk songs now in hip-hop style.)
\item The British abolished \textbf{Sati}, which orthodox Hindus resisted --- the \textbf{Dharma Sabha} (led by Radhakanta Deb) wanted to restore Sati, while the \textbf{Brahmo Samaj} (Raja Ram Mohan Roy, Western-educated) was against Sati.
\item \textbf{Wages for housewives/domestic work} are acceptable in Finland/Denmark/Sweden (feminist movements) but not to Indian culture, which portrays the wife as \textit{dharampatni} --- her dharma is to serve husband, child and elderly without return.
\item \textbf{Women in combat roles} in the army initially faced resistance.
\item The \textbf{French Revolution} was resisted by counter-enlightenment conservatives (Louis de Bonald, Joseph de Maistre, Auguste Comte); the \textbf{industrial revolution} was resisted by the \textbf{Luddites} (against technological development); electrification and computers were resisted by conservatives in India.
\end{itemize}
\item \textbf{Conservatives are the guardians of culture}, which is why the cultural system sits at the peak. Culture does not change overnight.
\end{itemize}
\end{CriticalPoint}

\begin{QuickRevision}
\begin{itemize}
\item \textbf{Cybernetic hierarchy of control} = theory of interrelation between action systems (subsystems).
\item \textbf{LIGA} order: L \& I high in \textbf{information/control}; G \& A high in \textbf{energy}. Information flows down; energy flows up.
\item Economy is highest in energy (like Marx's base). But economic/political change must be \textbf{acceptable to culture} (the regulatory authority).
\item Glocalisation (Ronald Robertson); Sati/Dharma Sabha vs Brahmo Samaj; Luddites; conservatives = guardians of culture.
\end{itemize}
\end{QuickRevision}

\sectiondivider

\subsection{Structural Differentiation \& Social Change}

Parsons gave relatively \textbf{less importance to social change}. He believed change is \textbf{evolutionary} (contrast: Marxists believe change is \textbf{revolutionary}).

\begin{KeyConcept}
\begin{itemize}
\item Whenever a change is introduced in any action system (A, G, I, L), it \textbf{develops strain}. The strain is then adjusted by the communication between LI and GA.
\item One line: \textbf{whenever any development takes place in society, it leads to strain; the other organs negotiate with the strain to restore equilibrium.}
\item Change is \textbf{evolutionary}, and when strain is negotiated with, a process called \textbf{structural differentiation} develops.
\end{itemize}
\end{KeyConcept}

\begin{KeyConcept}
\begin{itemize}
\item Parsons was much inspired by \textbf{Durkheim}. With industrialisation, people migrated to urban pockets, increasing \textbf{material density} (population size grew, resources limited) --- a strain. To negotiate this strain, \textbf{division of labour} (specialisation) evolved, which restored cooperation and equilibrium.
\item Durkheim did not call it ``structural differentiation''; Parsons calls the mitigation of strain by this process \textbf{structural differentiation} --- which can simply be read as \textbf{specialisation}.
\item Durkheim limited structural differentiation to the \textbf{division of labour}; Parsons extends it to the \textbf{division of the family into different specialised institutions}.
\end{itemize}
\end{KeyConcept}

\begin{ComparisonBox}
\begin{itemize}
\item \textbf{Durkheim:} specialisation of \textit{labour} $\rightarrow$ integration (division of labour).
\item \textbf{Parsons:} specialisation of \textit{organs/institutions} $\rightarrow$ the family divides into specialised institutions $\rightarrow$ interdependency $\rightarrow$ integration.
\end{itemize}
\end{ComparisonBox}

\paragraph{Structural differentiation of the family}

\begin{ExampleBox}
\textbf{Before industrialisation:}
\begin{itemize}
\item Only two organs ran society: the \textbf{classic extended family} (joint family) and the \textbf{church}.
\item The family performed multiple functions: caring for the elderly/dependents, economic production (cottage industry, farm work), nurturing/socialising the child, teaching skills, reproduction and love/sexual needs.
\item The \textbf{church} made the laws, gave education, appointed the ruler, gave blessings --- it was the regulatory authority.
\end{itemize}
\end{ExampleBox}

\begin{ExampleBox}
\textbf{After industrialisation (specialised institutions):}
\begin{itemize}
\item Elderly care $\rightarrow$ \textbf{care centres, HelpAge India}.
\item Nurturing/socialising the child $\rightarrow$ \textbf{kindergarten, play school, boarding school}.
\item Production $\rightarrow$ \textbf{industry} (family became a unit of \textit{consumption}, not production).
\item Teaching skills $\rightarrow$ \textbf{education/schooling system}.
\item Love/marriage $\rightarrow$ \textbf{dating apps, matrimony (shaadi.com)}.
\item Reproduction $\rightarrow$ \textbf{ART (artificial reproductive technologies), surrogacy}.
\item The \textbf{church} too lost its old functions (making laws, appointing the ruler, giving education) to the specialised \textbf{state} (with separation of church and state). Today the church's function is limited --- to religion alone, and specifically to \textbf{helping people manage the contradictions/anomalies of life} (e.g., a student who studies 12 hours a day but fails, while a 2-hour-a-day student qualifies; or an earthquake in which everyone dies except a two-year-old child).
\end{itemize}
\end{ExampleBox}

\begin{CriticalPoint}
\begin{itemize}
\item A common-sense view says the family is \textbf{declining} (joint $\rightarrow$ nuclear, families breaking). Parsons disagreed: it is not decline --- the family has \textbf{specialised into one specific domain} that no other institution can do: \textbf{primary socialisation} (roughly 0--4 years; feeding, walking, touching, smelling, first words). No institution has come for this.
\item So the family is left with \textbf{only primary socialisation} --- this is \textbf{structural differentiation of the family} (specialisation, not decline).
\item Health/education once given by the family are now given by the \textbf{state} --- right to education, national health services, Ayushman Bharat, National Health Mission, 35 kg free ration.
\end{itemize}
\end{CriticalPoint}

\begin{ExampleBox}
\begin{itemize}
\item Japan's rented relationship culture: A person will act like your girlfriend/boyfriend for one month, fully contractual and predetermined (e.g., \rupee 10,000 a month for agreed services). This is specialisation taken to an extreme.
\end{itemize}
\end{ExampleBox}

\begin{DataFact}
\begin{itemize}
\item \textbf{Goal attainment} is related to \textbf{polity/politics} --- politics sets the goal (e.g., Swachh Bharat 2019; making bureaucrats efficient via ``Mission Karmayogi''), and the \textbf{bureaucracy} executes it.
\item \textbf{Social action theory} was taken deeper by \textbf{Niklas Luhmann} (not in this syllabus).
\item \textbf{Mode of control} = a mechanism of social control, part of normative constraint. The school teaches obedience and conformity through the \textbf{hidden curriculum} (good morning, discipline, clean appearance) --- a mechanism of social control.
\item \textbf{Energy} simply means sudden change (e.g., industrialisation = high energy).
\item \textbf{Jean Piaget's} cognitive-development theory deals with the child's stages, \textbf{not} society's development. Parsons was \textbf{not} influenced by Piaget; he was influenced by \textbf{Sigmund Freud} (defence mechanism --- society's organs ``cool down'' the energy of change, like the ego's defence mechanism).
\item Parsons was originally a \textbf{biologist}, not a sociologist (he had also studied economics). Durkheim was a philosopher; Weber became a sociologist only after 1905; Marx was never a sociologist. Merton, by contrast, was academically \textit{only} a sociologist.
\item The \textbf{goal of tribal society} is the same as any society's --- to sustain development and survival (hunting, gathering, ensuring raw food is available).
\end{itemize}
\end{DataFact}

\begin{QuickRevision}
\begin{itemize}
\item Parsons: change is \textbf{evolutionary} (vs Marx's revolutionary). Change $\rightarrow$ strain $\rightarrow$ structural differentiation $\rightarrow$ equilibrium.
\item \textbf{Structural differentiation} = specialisation; borrowed from Durkheim's division of labour, but extended to division of the family into institutions.
\item Family specialised into \textbf{primary socialisation} only $\rightarrow$ not decline but specialisation.
\item Parsons = biologist influenced by Freud; change = high energy; hidden curriculum = social control.
\end{itemize}
\end{QuickRevision}

\sectiondivider

\section{Talcott Parsons --- Social Change}

\subsection{Revision: AGIL Schema \& Organ Functions}

This lecture opens with a quick revision of the \textbf{AGIL schema} --- the four needs of society, each fulfilled by a specific organ, each contributing to the maintenance of the system:

\begin{center}
\rowcolors{2}{white}{Tint}
\begin{tabularx}{\textwidth}{p{3.2cm} p{4cm} Y}
\rowcolor{AccentDark}
\thead{Need} & \thead{Organ} & \thead{How it maintains the system} \\
\midrule
\textbf{Adaptation} & Economy & Producing and distributing goods and resources \\
\textbf{Goal attainment} & Polity & Establishing goals and mobilising resources to fulfil them \\
\textbf{Integration} & Law, religion, education & Restoring equilibrium; resolving conflict; bringing people together \\
\textbf{Latency} & Education, family, religion & Managing tension \\
\bottomrule
\end{tabularx}
\end{center}

\begin{ExampleBox}
\begin{itemize}
\item \textbf{Economy (polity's goal example):} ``Make in India'' --- to make India a manufacturing hub, resources must be mobilised; e.g., an FDI clause requiring a certain percentage of inputs to be sourced from Indian soil gives employment and pushes the domestic economy (foreign investment integrated with, not at the cost of, the domestic economy).
\item \textbf{Law:} a family dispute goes to a \textbf{Lok Adalat} (court for personal/family/domestic issues) and is resolved through arbitration and mediation, restoring equilibrium.
\item \textbf{Religion:} during crisis (death of a loved one), people gather, express grief, conduct rituals, pray for the soul --- integrating people.
\item \textbf{Education (latency):} tension management.
\end{itemize}
\end{ExampleBox}

\begin{QuickRevision}
\begin{itemize}
\item AGIL: A$\rightarrow$economy (production/distribution), G$\rightarrow$polity (goal-setting + resource mobilisation), I$\rightarrow$law/religion/education (conflict resolution), L$\rightarrow$education/family/religion (tension).
\item Make in India = goal attainment; Lok Adalat = law's integrative role.
\end{itemize}
\end{QuickRevision}

\sectiondivider

\subsection{The Problem of Explaining Social Change}

The central question: \textbf{if everything is so rosy, and every conflict is resolved, how do we explain social change?} Parsons does explain social change --- with the same theory.

\begin{CriticalPoint}
\begin{itemize}
\item Parsons is a scholar who explains social change \textbf{with this same equilibrium theory}. The revision of cybernetic hierarchy is necessary because it tells us (a) the relationship between the four action systems, and (b) \textbf{why a system is maintained the way it is}.
\end{itemize}
\end{CriticalPoint}

\begin{QuickRevision}
\begin{itemize}
\item If every issue is resolved by the system, change would never come --- so how does Parsons explain change? Through cybernetic hierarchy + the concept of strain.
\end{itemize}
\end{QuickRevision}

\sectiondivider

\subsection{Cybernetic Hierarchy \& the Energy Balance}

\begin{KeyConcept}
\begin{itemize}
\item Two systems (\textbf{L, I}) are high in \textbf{information}; two (\textbf{G, A}) are high in \textbf{energy}.
\item The \textbf{input from the economy gives energy} to the system, and the \textbf{output from the cultural system gives resistance} to that input. This \textbf{input--input--output balance} maintains the system.
\item \textbf{Excess of energy} (too much change in the economy) creates trouble; \textbf{deficiency of energy} (economic depression) also creates trouble.
\end{itemize}
\end{KeyConcept}

\begin{DataFact}
\begin{itemize}
\item 3--4\% inflation is healthy --- it shows a demand-driven economy, consumerism, purchasing power, productivity.
\item Beyond 8--10\% inflation becomes a problem --- it shows lack of supply.
\item Hyperinflation is not okay; depression (deficiency) is also not okay. \textbf{Excess of anything is not okay} --- and excess can change the system.
\end{itemize}
\end{DataFact}

\begin{QuickRevision}
\begin{itemize}
\item L/I = high information; G/A = high energy. Economy (energy input) $\leftrightarrow$ culture (resistance/output) balance maintains the system.
\item Excess energy (hyperinflation) or deficiency (depression) = trouble = change.
\end{itemize}
\end{QuickRevision}

\sectiondivider

\subsection{Factors Causing Strain (Five Factors)}

Strain develops because of \textbf{excess of energy or deficiency of energy}. What exactly causes excess/deficiency? Five factors:

\begin{DataFact}
\begin{enumerate}
\item \textbf{Changes in the demographic character of the population} --- population explosion, sudden waves of migration, exodus/mass migration, refugee crisis, rising infant mortality, racial intermixture (intermarriage), and changes in mortality and fertility rates. (Note: high infant mortality at a certain level is ``fine'' --- beyond that it becomes a \textbf{public issue}, not personal trouble.)
\item \textbf{Changes in the physical environment} --- exhaustion of physical resources (soil, water, weather conditions).
\item \textbf{Changes in population resulting from increased productivity of food and availability of resources} --- increased food productivity $\rightarrow$ population rise $\rightarrow$ migration $\rightarrow$ load on land $\rightarrow$ conflict between indigenous and ``alien'' communities $\rightarrow$ sons-of-the-soil and secessionist movements.
\item \textbf{Changes in technology and application of scientific knowledge} --- technological advancement creates strain.
\item \textbf{Development of new cultural configuration} --- new religious ideas, or integration of religious values with science.
\end{enumerate}
\end{DataFact}

\begin{ExampleBox}
\textbf{Physical environment (Punjab \& Haryana):}
\begin{itemize}
\item Punjab/Haryana faced food shortage due to exhaustion of physical resources (bad health, saline soil, groundwater shortage, excess irrigation/fertilisers, famine).
\item The \textbf{Green Revolution} (high-yielding-variety seeds) helped them cope.
\item \textbf{Negative repercussions:} female child/infant mortality rose (female children seen as burden; people could not feed themselves); migrant labourers were underpaid and exploited; \textbf{fundamentalism emerged (Khalistan movement)}; the demand for a separate state and the Punjab--Haryana divide grew out of the \textbf{Sutlej Canal} water-sharing dispute.
\end{itemize}
\end{ExampleBox}

\begin{ExampleBox}
\textbf{Technology:}
\begin{itemize}
\item The \textbf{first/second/third-world division} and global inequality are the result of scientific/technological development --- the third world depends on the first world (defence equipment from Israel, agricultural inputs), which is why ``Make in India''/``vocal for local'' is pushed.
\item \textbf{Sex-determination tests} lead to female foeticide.
\item \textbf{Artificial intelligence} --- no one knows whether it will create or destroy employment. Some studies (e.g., \textbf{Deloitte}) say employment will rise in certain (developed-country) sectors; in India's huge informal sector (which runs on cheap, unsecured labour), AI may be the biggest source of unemployment. Informal-sector workers (Uber/Swiggy drivers) cannot even form trade unions --- a concept Michael Burawoy explains through ``work turned into a game'' (each worker has an individual target and bonus, so work feels like a game, not exploitation --- this minimises the idea of trade union).
\end{itemize}
\end{ExampleBox}

\begin{ExampleBox}
\textbf{New cultural configuration:}
\begin{itemize}
\item \textbf{Martin Luther} challenged the Catholic Church $\rightarrow$ he was called a heretic and excommunicated. He motivated German peasants, who were looted through \textbf{tithe} (first 10\%, then 20\%, justified by religion --- ``advance to the next level to go to heaven''). When Luther translated the Bible into English, peasants realised it was not written there, and a \textbf{peasant revolt} (mass execution) followed.
\item \textbf{Newton, Galileo, Copernicus, Kepler} --- declared heretics for discoveries against religious ideas.
\item \textbf{Buddha} challenged Hinduism $\rightarrow$ Brahminical priests defamed him.
\end{itemize}
\end{ExampleBox}

\begin{CriticalPoint}
\begin{itemize}
\item Change can be initiated by \textbf{both economy and superstructure} --- strain can come from excess/deficiency of energy \textit{or} information, i.e., from culture, economy, or anywhere (e.g., Buddha's challenge was cultural change).
\end{itemize}
\end{CriticalPoint}

\begin{QuickRevision}
\begin{itemize}
\item Five strain factors: (1) demographic character, (2) physical environment, (3) population from food productivity, (4) technology/science, (5) new cultural configuration.
\item Green Revolution: solved food shortage but caused female foeticide, exploitation, Khalistan movement, Punjab--Haryana divide.
\item Change can start from economy \textit{or} culture (energy or information).
\end{itemize}
\end{QuickRevision}

\sectiondivider

\subsection{The Four Processes of Change}

When strain develops, it leads to \textbf{four processes} (memorise these --- they are the keywords):
\begin{enumerate}
\item \textbf{Structural differentiation}
\item \textbf{Integration}
\item \textbf{Adaptive upgradation}
\item \textbf{Generalisation of values}
\end{enumerate}
All four are \textbf{evolutionary}, not revolutionary.

\begin{CriticalPoint}
\begin{itemize}
\item ``Revolution'' is an ideal type (a Weberian concept). Parsons' account has nothing revolutionary in it --- change is always evolutionary.
\end{itemize}
\end{CriticalPoint}

\paragraph{1. Structural differentiation}
When strain develops, organs \textbf{specialise} to cope --- they lose some previous functions to specialised organs.

\begin{ExampleBox}
\begin{itemize}
\item India came up with specialised institutions: quarantine centres, COVID care centres; hotels, schools and railway stations were \textbf{transformed into hospitals}; hospital beds increased; new hospitals institutionalised. (Nothing is objective reality --- time and space change the seemingly objective reality.)
\item This is the \textbf{adaptive process} India underwent.
\end{itemize}
\end{ExampleBox}

\paragraph{2. Integration}
Specialised organs automatically create \textbf{interdependency}, hence integration.

\begin{KeyConcept}
\begin{itemize}
\item \textbf{Adam Smith:} the function of division of labour is \textbf{productivity} (Wealth of Nations).
\item \textbf{Durkheim:} the function of specialisation is \textbf{integration} (organic solidarity).
\item \textbf{Parsons combines both:} the function of specialisation is \textbf{both integration and productivity}. The more specialised institutions a society has, the more developed it is.
\end{itemize}
\end{KeyConcept}

\paragraph{3. Adaptive upgradation}

\begin{ExampleBox}
\begin{itemize}
\item During COVID, India gave \textbf{honorarium to medical warriors} and increased the wages of \textbf{ASHA workers, auxiliary nurse midwives (ANMs)}, etc.\ --- adaptive upgradation.
\end{itemize}
\end{ExampleBox}

\paragraph{4. Generalisation of values}
To integrate all the specialised organs, members must \textbf{share certain values} (shared consensus).

\begin{ExampleBox}
\begin{itemize}
\item An engineer and a policeman are each specialised in different tasks; they integrate because each depends on the other, \textbf{provided} they share values (democracy, equal opportunity). If a person believes the police/army only support the rich, integration fails.
\item As an American, one must believe in \textbf{American democratic values} (Americanism, democracy, capitalism) --- a Marxist/left-winger has ``no space'' in America.
\item In USA movies, whenever anything happens, one calls \textbf{9-1-1} --- the shared value of trusting the system (a Naxal would \textit{not} call the police).
\end{itemize}
\end{ExampleBox}

\begin{DataFact}
\textbf{The diagram (Parsons' Theory of Social Change):}
\begin{itemize}
\item Heading: \textbf{``Parsons' Theory of Social Change''} (also called \textbf{``Functional Theory of Social Change''}), built into the \textbf{Theory of Cybernetic Hierarchy of Control}.
\item Box 1: Excess of energy or information, or insufficient supply of energy or information, \textbf{causes strain} in the system.
\item Box 2: Strain leads to \textbf{structural differentiation, integration, adaptive upgradation, generalisation of values}.
\item Box 3: These four processes lead to the establishment of a \textbf{new equilibrium}.
\end{itemize}
\end{DataFact}

\begin{DataFact}
\textbf{America as the model:}
\begin{itemize}
\item Asked which country is most developed, Parsons would answer \textbf{America} --- not because he is American, but because America has the \textbf{highest adaptive capacity}: it adapts to extreme changes and creates specialised institutions in no time.
\item Example: \textbf{Moderna} was the first country-company to develop the COVID vaccine, showing the highest adaptive capacity to the COVID strain.
\end{itemize}
\end{DataFact}

\begin{QuickRevision}
\begin{itemize}
\item Strain $\rightarrow$ four processes: \textbf{structural differentiation, integration, adaptive upgradation, generalisation of values} (all evolutionary).
\item Revolution = ideal type (nothing revolutionary here).
\item Specialisation $\rightarrow$ interdependency $\rightarrow$ integration (Smith's productivity + Durkheim's integration combined).
\item America = highest adaptive capacity (Moderna vaccine); generalisation of values = shared consensus (9-1-1).
\end{itemize}
\end{QuickRevision}

\sectiondivider

\section{Critique of Parsons \& Answer-Writing}

\subsection{Answer-Writing: ``Power and Authority Go Together''}

The lecture opens with a 10-mark question: \textbf{``Power and authority go together. Examine.''}

\begin{PYQLink}
\begin{itemize}
\item The demand of the question is \textbf{not} just ``power'' and ``authority'' --- it is \textbf{``go together''} (the relation). If you explain only the two terms and not the relation, you escape the demand and lose marks.
\end{itemize}
\end{PYQLink}

\paragraph{The approach (answer structure)}
\begin{enumerate}
\item \textbf{Identify keywords} --- power, authority (and the relation ``go together'').
\item \textbf{Define power:} the \textit{probability of a man (or a number of men) to exercise his will even against the resistance of others}.
\item \textbf{Define authority:} the \textit{probability of obedience to a command}.
\begin{itemize}
\item These two definitions fetch about 2--2.5 marks out of 10.
\end{itemize}
\item \textbf{Explain the relation:} power, when it gets legitimacy, becomes authority --- \textbf{Power + Legitimacy = Authority (P + L = A)}.
\item \textbf{Explain the three bases of legitimacy} (and thus three forms of authority):
\begin{itemize}
\item \textbf{Traditional} authority
\item \textbf{Charismatic} authority
\item \textbf{Legal-rational} authority
\item (Give one example for each.)
\end{itemize}
\item \textbf{Conclusion:} the larger trend of societies is towards \textbf{rationalisation} --- power is increasingly legitimised on \textbf{legal-rational} principles, less and less on tradition.
\end{enumerate}

\begin{ExampleBox}
\begin{itemize}
\item A man demands \rupee 2,000 from you for not wearing a mask (power over you), and you resist. The moment he shows his \textbf{ID card} (legitimacy --- he is a policeman), you obey. The card gives legitimacy; power + legitimacy = authority.
\end{itemize}
\end{ExampleBox}

\begin{CriticalPoint}
\begin{itemize}
\item Keep the answer to $\sim$150 words. You \textit{may} add extra things (social action, ideal type, fashion) but you must always \textbf{return to the demand of the question} --- do not divert.
\end{itemize}
\end{CriticalPoint}

\begin{QuickRevision}
\begin{itemize}
\item Power = probability of exercising will against resistance; Authority = probability of obedience to command.
\item \textbf{Power + Legitimacy = Authority.}
\item Three legitimacies $\rightarrow$ three authorities: traditional, charismatic, legal-rational.
\item Conclusion: rationalisation trend $\rightarrow$ legal-rational legitimacy rising.
\end{itemize}
\end{QuickRevision}

\sectiondivider

\subsection{Structural Differentiation of Caste}

After the family-and-church example, Parsons' idea of structural differentiation is applied to \textbf{caste in India}.

\begin{DataFact}
\textbf{Functions of caste in ancient times:}
\begin{itemize}
\item Caste defined your occupation, status, class, power, marriage alliance, socialisation and education. Education was highly casteist --- upper castes learned their occupations, lower castes learned menial tasks, so lower castes reproduced their own occupations (e.g., sanitation workers' children did the same work).
\end{itemize}
\end{DataFact}

\begin{ExampleBox}
\textbf{Caste today (structural differentiation):}
\begin{itemize}
\item Today occupation is defined by \textbf{individual merit, skill and education} (not caste); status comes from occupation/work ethics (an IAS officer has status, not because of caste); power/status no longer flow from caste.
\item Marriage is still largely caste-based (one domain caste retains), but occupation/status/education are no longer caste-defined.
\item So caste has not ``declined'' --- it has \textbf{specialised into one or two domains}, chiefly \textbf{democratic politics} (reservation, dominant caste, caste-based electoral politics). This is \textbf{structural differentiation of caste}.
\end{itemize}
\end{ExampleBox}

\begin{CriticalPoint}
\textbf{Indian modernity:}
\begin{itemize}
\item Indian society has \textbf{both} pattern variable A and B. We have democracy (pattern variable B), but you cannot yield maximum fruits from democracy without using caste in electoral politics (pattern variable A). Therefore Indian modernity is \textbf{an Indian modernity} --- not a European one we emulate, but a unique mix of both worlds.
\end{itemize}
\end{CriticalPoint}

\begin{ComparisonBox}
\textbf{Integration between caste and democracy:}
\begin{itemize}
\item Caste and democracy are two ``opposite'' institutions, yet they integrate: \textbf{democracy needs caste} (caste is a vote bank) and \textbf{caste needs democracy} (reservation gives lower castes positions of power).
\item Famous phrase: \textbf{``In India, we don't caste our votes; we vote our caste.''}
\item \textbf{Paradox:} affirmative action (reservation) gives lower castes equal opportunity, but opportunities are partly given \textbf{on the basis of caste} --- a paradox.
\end{itemize}
\end{ComparisonBox}

\begin{CriticalPoint}
\textbf{Generalisation of values:}
\begin{itemize}
\item Parsons believed \textbf{America} has fully internalised democratic virtues (no role of religion/race/ascriptive features in politics or employment --- only achievement).
\item We \textbf{cannot be so sure about India}: religion, caste and patriarchy still govern politics, so India may not have internalised democratic virtues.
\item \textbf{School is a focal agency / bridge between family and society} --- it spreads values and prepares the child for wider society (critical thinking, debate). American schools (Harvard, Cambridge = ``schools'') are the best.
\end{itemize}
\end{CriticalPoint}

\begin{DataFact}
\textbf{American religion (Robert Bellah):}
\begin{itemize}
\item In America there are \textbf{private/individual religions} (individualistic beliefs of modern society), but also one overarching \textbf{civic religion --- ``Americanism''} (nationalism), which integrates and cuts across all the private religions and unifies them. This is a concept from \textbf{Robert Bellah} (two kinds of religion).
\end{itemize}
\end{DataFact}

\begin{ExampleBox}
\textbf{America's adaptive capacity (again):}
\begin{itemize}
\item Whatever problem America faces (immigration under Trump's ``America first'', hurricanes/polar vortex, etc.), since America has the highest adaptive capacity, equilibrium is restored. In the 1930 \textbf{Great Depression}, America produced \textbf{John Maynard Keynes} (monetary/fiscal policy) and \textbf{Roosevelt} (who reposed faith in Keynesian economics) --- equilibrium restored.
\end{itemize}
\end{ExampleBox}

\begin{QuickRevision}
\begin{itemize}
\item Caste lost most functions (occupation/status/education) $\rightarrow$ specialised into \textbf{democratic politics} = structural differentiation of caste.
\item Indian modernity = mix of pattern variables A \& B (``we don't caste our votes, we vote our caste'').
\item School = \textbf{focal agency/bridge} between family and society.
\item America = highest adaptive capacity (Keynes, Roosevelt in Great Depression).
\end{itemize}
\end{QuickRevision}

\sectiondivider

\subsection{Moving Equilibrium \& Evolutionary Universals}

The four processes lead to a \textbf{new equilibrium}. Every society has an inbuilt capacity to adjust --- Parsons calls this \textbf{homeostasis}.

\begin{KeyConcept}
\textbf{Moving equilibrium:}
\begin{itemize}
\item Before strain, there is an old equilibrium; after adjustment, a new equilibrium. Then new strains arise, and again a new equilibrium. So society is \textbf{a story of moving equilibrium}.
\item ``Moving'' and ``equilibrium'' are opposites --- Parsons uses this paradoxical pair to mean society \textbf{changes continuously, not drastically}. It is also called \textbf{dynamic equilibrium}.
\item This is \textbf{change in continuity} --- gradual, evolutionary, continuous change, not revolutionary.
\end{itemize}
\end{KeyConcept}

\begin{KeyConcept}
\textbf{Two levels of adjustment:}
\begin{itemize}
\item Readjustment occurs at \textbf{two levels}:
\begin{enumerate}
\item \textbf{Change within the system} (continuous change inside a stage).
\item \textbf{Change of the system} (moving from one stage to another).
\end{enumerate}
\item (Note: ``change of the system'' is deliberately kept light --- the syllabus focuses on change within.)
\end{itemize}
\end{KeyConcept}

\begin{DataFact}
\textbf{Evolutionary universals (stages):}
\begin{itemize}
\item Society's evolution is divided into stages --- like an individual's continuous growth divided into \textbf{infant $\rightarrow$ teenage $\rightarrow$ adult}.
\item Parsons' stages: \textbf{primitive $\rightarrow$ intermediate $\rightarrow$ modern} (infant $\rightarrow$ teenage $\rightarrow$ adult).
\item These stages are called \textbf{evolutionary universals} --- every society (American, Indian, Chinese) universally passes through all three, in a linear fashion.
\item Parsons believed America is the ``adult'' (highest adaptive capacity); India and China are ``intermediate'' (second world).
\end{itemize}
\end{DataFact}

\begin{KeyConcept}
\textbf{Modernization theory:}
\begin{itemize}
\item From Parsons' writings came \textbf{modernisation theory}: how America became modern, and what others should learn. America became the ``guardian of development'' --- hence institutions like the \textbf{IMF, World Bank, Marshall Plan}, and soft/hard loans, with conditions attached (``if you want help with balance-of-payments crisis, listen to us'').
\item Modernising meant \textbf{becoming like America} --- development became synonymous with America, westernisation with Americanisation.
\item (Note: America \textit{was} once primitive too --- but its primitive history (Red Indians) is localised to tribes and studied only by anthropologists.)
\end{itemize}
\end{KeyConcept}

\begin{CriticalPoint}
\textbf{Scholars evolve:}
\begin{itemize}
\item Scholars evolve: \textbf{young Marx vs old Marx} (the \textbf{epistemological break} --- \textit{Economic and Philosophical Manuscripts} on alienation vs later theory of revolution). Gandhi also changed his views over time. Parsons too has a ``young'' and ``old'' version, but those details are beyond this syllabus.
\item In the longer run, readjustment takes an \textbf{evolutionary form}; since strain factors are constantly at work, readjustment is continuous --- hence society is in \textbf{moving equilibrium}.
\end{itemize}
\end{CriticalPoint}

\begin{ExampleBox}
\textbf{The relevance of Parsons to everyday life (population control):}
\begin{itemize}
\item When a state government (e.g., UP, Assam) brings a \textbf{population-control policy} in a coercive/Malthusian fashion (forcing sterilisation, as in the mass campaign during the 1975 Emergency), it \textbf{creates strain} --- one of the factors causing strain is change in population.
\item To restore equilibrium, \textbf{civil society, NGOs, pressure groups, editorial/newspaper/media} mobilise to oppose the policy, and it may never be enacted. There is also a \textbf{Population Stabilisation Fund} --- a move towards equilibrium.
\item The better (and more effective) approach is \textbf{bottom-up}, not top-down: \textbf{educating women} (literate women keep population in check) --- hence coercive population policies do not work. (These state governments are, in this sense, ``new Malthusians'' --- Malthus was also a positivist scholar.)
\end{itemize}
\end{ExampleBox}

\begin{QuickRevision}
\begin{itemize}
\item Four processes $\rightarrow$ \textbf{new equilibrium}; society = \textbf{moving/dynamic equilibrium} = change in continuity.
\item Two adjustment levels: \textbf{change within} vs \textbf{change of} the system.
\item \textbf{Evolutionary universals:} primitive $\rightarrow$ intermediate $\rightarrow$ modern (universal, linear).
\item Modernisation theory: development = becoming America (IMF, World Bank, Marshall Plan).
\item Coercive population control $\rightarrow$ strain $\rightarrow$ civil society/media restore equilibrium; educate women (bottom-up).
\end{itemize}
\end{QuickRevision}

\sectiondivider

\subsection{Cognitive Consonance: Pattern Variables \& Social Systems}

This is the ``patchwork'' that connects the two things studied: \textbf{pattern variables} and the \textbf{social system} --- their \textbf{cognitive consonance} (intellectual relation; they ``go hand in hand'').

\begin{KeyConcept}
\begin{itemize}
\item Any entity is a social system if it has the \textbf{four organs (AGIL)} catering to the four needs --- India, a university, even a bureaucracy.
\item Every social system runs on a \textbf{specific set of pattern variables}.
\end{itemize}
\end{KeyConcept}

\begin{ComparisonBox}
\textbf{Bureaucracy (pattern variable B):}
\begin{itemize}
\item A bureaucracy is a social system. It is most \textbf{efficient} when it runs on pattern variable B: \textbf{universalism} (merit-based recruitment --- ``give the exam, even if you are my nephew''), \textbf{achievement orientation}, \textbf{self-orientation}, and \textbf{affective neutrality}.
\item \textbf{Affective neutrality:} an interview board is cool, objective and instrumental towards the candidate --- not emotional or empathetic.
\end{itemize}
\end{ComparisonBox}

\begin{ComparisonBox}
\textbf{Family (pattern variable A):}
\begin{itemize}
\item The family is also a social system with its own AGIL (adaptation: mother as household guardian, father as breadwinner; goal attainment: parents' goals and resource mobilisation; integration: resolving sibling conflicts; latency: motivating and diffusing tension).
\item But the family \textbf{cannot} run on pattern variable B (bureaucratic principles) --- it would break down. The family must run on \textbf{pattern variable A} (expressive).
\item \textbf{Gendered role expectations:} mother performs the \textbf{expressive role}, father the \textbf{instrumental role}. Pattern variable A is expressive; B is instrumental (analogous to Weber's formal rationality --- the principles for efficient bureaucracy).
\end{itemize}
\end{ComparisonBox}

\begin{QuickRevision}
\begin{itemize}
\item \textbf{Cognitive consonance} = pattern variables and social systems go hand in hand.
\item Bureaucracy $\rightarrow$ efficient on pattern variable B (universalism, achievement, affective neutrality).
\item Family $\rightarrow$ runs on pattern variable A (expressive; mother expressive, father instrumental). Running on B breaks it.
\end{itemize}
\end{QuickRevision}

\sectiondivider

\subsection{Feminism \& the Decline of the Family}

\begin{CriticalPoint}
\begin{itemize}
\item Family survives on the two gendered role expectations (mother expressive, father instrumental). With \textbf{feminism}, the mother also wants the instrumental role (occupation) --- now both go instrumental, and the child's care goes to \textbf{crèche/day-care centres}.
\item In America, the family is \textbf{declining} (in the doldrums) only because of feminism.
\item \textbf{Why family declines under capitalism:} to be an efficient worker you must be mobile (ready to relocate to Bangalore/Chennai/Silicon Valley), which weakens family; you find a new partner and start a \textbf{live-in relationship} (to save money); then job offers create \textbf{long-distance relationships} (not under a common roof); finally \textbf{LAT --- ``living apart together''} (married but staying apart for work). (DINK --- ``double income, no children'' --- is more slang.)
\end{itemize}
\end{CriticalPoint}

\begin{ExampleBox}
\begin{itemize}
\item Durkheim is always concerned with \textbf{cohesion/integration}. Feminism \textbf{disintegrates} (does not integrate), so Durkheim would never like feminism.
\end{itemize}
\end{ExampleBox}

\begin{QuickRevision}
\begin{itemize}
\item Family needs gendered role expectations (expressive mother + instrumental father). Feminism $\rightarrow$ both instrumental $\rightarrow$ family declines.
\item Decline chain: family $\rightarrow$ live-in $\rightarrow$ long-distance $\rightarrow$ \textbf{LAT (living apart together)}.
\item Durkheim would reject feminism (it disintegrates).
\end{itemize}
\end{QuickRevision}

\sectiondivider

\subsection{Critique of Parsons}

\begin{CriticalPoint}
\textbf{The core weakness:}
\begin{itemize}
\item Parsons explains school, family, and all organs functioning perfectly $\rightarrow$ a great society. He uses \textbf{inductive reasoning} (like a positivist), believing \textbf{every} society follows the primitive $\rightarrow$ intermediate $\rightarrow$ modern path. His theory cannot explain recurrent real-world problems.
\end{itemize}
\end{CriticalPoint}

\paragraph{Critique 1 --- George Huaco (structural functionalism as a political program)}
\begin{CriticalPoint}
\begin{itemize}
\item \textbf{George Huaco} says structural functionalism is a \textbf{political program} --- developed (and, he claims, funded by the American state) to \textbf{legitimise American hegemony} over other countries. It portrays America as the most effective nation at dealing with strain, and others as less effective, so they should follow America.
\end{itemize}
\end{CriticalPoint}

\paragraph{Critique 2 --- Ralf Dahrendorf (conflict theorist)}
\begin{CriticalPoint}
\begin{itemize}
\item \textbf{Dahrendorf} (a Marxist/conflict theorist) says Parsons is a \textbf{``veiled status quoist''} --- he talks about change but always returns to equilibrium (status quo). He ignores \textbf{conflict between sections} (man--woman, rich--poor, Hindu--Muslim) and talks only about abstract models, not empirical reality.
\item When a conflict theorist reads the newspaper, he sees only \textbf{conflict}, not ``maintenance of system'' --- so \textbf{conflict is the essence of society}, not consensus. Parsons does not deal with real conflict over power between groups.
\end{itemize}
\end{CriticalPoint}

\paragraph{Critique 3 --- Jeffrey Alexander (neo-functionalism)}
\begin{CriticalPoint}
\begin{itemize}
\item \textbf{Jeffrey Alexander} (who gave \textbf{neo-functionalism} --- a theory he himself later said not to read) argues Parsons \textbf{cannot explain real social issues}: rise of fascist leaders, communal conflict, sexual harassment of women (private and public sphere), divorce, crime against women, unemployment, and school dropouts.
\item Parsons would only call these ``strain'' and say the political system will fix them --- a patchwork, not an explanation.
\end{itemize}
\end{CriticalPoint}

\begin{DataFact}
\textbf{Related points (from other lectures):}
\begin{itemize}
\item \textbf{C. Wright Mills} (who introduced Marxism into America) argued structural functionalism is \textbf{teleological} (explaining effect in the name of cause --- ``a chair is to sit on'' explains its \textit{use}, not the chair) and \textbf{tautological} (saying the same thing in many words). \textbf{Alvin Gouldner} wrote \textit{The Coming Crisis of Western Sociology} --- the decline of structural functionalism (with the rise of Marxism) was the ``crisis of Western sociology.''
\end{itemize}
\end{DataFact}

\begin{CriticalPoint}
\begin{itemize}
\item Parsons speaks like a positivist (inductive reasoning; all societies evolve primitive $\rightarrow$ intermediate $\rightarrow$ modern; equilibrium and gradual evolution). He was inspired by three streams --- \textbf{positivism, utilitarianism, and idealism} --- and you see elements of all three in his concepts.
\end{itemize}
\end{CriticalPoint}

\begin{QuickRevision}
\begin{itemize}
\item Three critiques: \textbf{Huaco} (political program for American hegemony), \textbf{Dahrendorf} (veiled status quoist, ignores conflict), \textbf{Alexander} (can't explain fascism, communalism, sexual harassment, unemployment, dropouts).
\item Plus \textbf{C. Wright Mills} (teleological/tautological) and \textbf{Alvin Gouldner} (\textit{Coming Crisis of Western Sociology}).
\item Parsons = positivist (positivism + utilitarianism + idealism).
\end{itemize}
\end{QuickRevision}

\sectiondivider

\subsection{Test Question Discussion (Answer-Writing Toolkit)}

\paragraph{Q. ``Sociology developed as a conservative reaction to modernity. Do you agree?''}
\begin{enumerate}
\item \textbf{Define modernity} (1 mark) --- modernity is the central issue.
\item \textbf{Trace the evolution of modernity} to revolutionary social changes in Europe --- scientific revolution, intellectual revolution, French Revolution, industrial revolution (half mark).
\item \textbf{List the features/changes} (+2 marks), e.g.:
\begin{itemize}
\item Patriarchy $\rightarrow$ gender equality (\textbf{Mary Wollstonecraft})
\item Rule of religion $\rightarrow$ secularism (\textbf{Voltaire})
\item Absolute monarchy $\rightarrow$ constitutional monarchy (\textbf{Locke, Rousseau})
\item God-centric $\rightarrow$ man-centric world (\textbf{humanism})
\item Joint family $\rightarrow$ nuclear family
\end{itemize}
\item \textbf{Conservative reaction} (counter-revolutionary phase) --- French intellectuals \textbf{de Bonald, de Maistre, Comte, Durkheim}: de Bonald argued traditional institutions maintained social order and were functional; Comte believed social laws govern the world and can be uncovered; Durkheim proposed social reforms to fix pathologies. $\rightarrow$ This led to the rise of \textbf{structural functionalism}.
\item \textbf{Radical reaction} --- social change also invoked a radical reaction, later called \textbf{Marxism} (revolutionary, unlike structural functionalism).
\item \textbf{Conclusion} --- these unique perspectives around modernity led to the emergence of sociology.
\end{enumerate}

\paragraph{Q. ``Ethnomethodology is a way of doing phenomenological sociology.''}
\begin{enumerate}
\item \textbf{Define ethnomethodology} --- the study of the \textit{methods used by members of society to make sense of everyday life}.
\item \textbf{Phenomenology} --- a response to positivism; the social world is a \textbf{web of meanings} with no objective reality or laws governing actors.
\item \textbf{Ethnomethodology as a branch of phenomenology} --- developed by \textbf{Garfinkel}; phenomenology is the theory, ethnomethodology is the laboratory/practice.
\item \textbf{Three methods} (with examples):
\begin{itemize}
\item \textbf{Breaching experiment} --- breaching social order to uncover the chaos beneath (e.g., playing the national anthem in a shopping mall; a bearded man wearing a sari at JNU = resistance to the heteronormative gender order).
\item \textbf{Documentary method} --- repetitive behaviour makes us believe reality is objective (girls in sari, boys in shirt-pants $\rightarrow$ constructed norms, not objective reality).
\item \textbf{Indexicality} --- meaning depends on context (a boy in a sari in a theatre = a play; outside the theatre = shock).
\end{itemize}
\item \textbf{Conclusion} --- ethnomethodology is doing phenomenology. (Scholars referenced: \textbf{Yoginder Singh}, \textit{Modernization of Indian Tradition}.)
\end{enumerate}

\paragraph{Q. ``Contextualize sociology in the debate between science and common sense.''}
\begin{enumerate}
\item \textbf{Differentiate common sense vs science} (a table).
\item \textbf{Locate sociology} --- sociology fits \textbf{both}: phenomenology/symbolic interactionism/ethnomethodology are close to \textbf{common sense}; positivism is close to \textbf{science}; \textbf{interpretative sociology} bridges all three (science of social action).
\item \textbf{Conclusion} --- sociology is a \textbf{multi-paradigmatic science} (no single dominant paradigm, unlike natural science). \textbf{Lyotard} adds that science itself is a meta-narrative (not objective today).
\end{enumerate}

\paragraph{Q. ``Higher unemployment rates are the source of higher crime rates. Which theoretical strand would support this? Is it scientific?''}
\begin{enumerate}
\item \textbf{Define ``scientific''} --- per \textbf{Popper}, a statement is scientific only if it can be \textbf{proven wrong} (falsifiability).
\item \textbf{Positivism} --- the proposition is positivist (quantitative methodology; social facts govern individuals; causal relation; social engineering). Positivist examples: Comte's law of three stages; Durkheim's law on suicide (industrialisation $\rightarrow$ higher suicide rate).
\item \textbf{Non-positivist critique} --- it is a \textbf{monocausal} explanation of a multi-causal world; the relation can be inverted (crime can also cause unemployment); it is contextual (USA has low unemployment but higher crime than India; urban India has higher crime than rural); the \textbf{social location of the researcher} (Alvin Gouldner) influences definitions of crime/unemployment (feminist/Marxist researchers); Marxist critique --- it conceals bourgeois crime (``profit is theft''), blames the working class, and ignores the industrial reserve army.
\item \textbf{Conclusion} --- because the statement \textit{can} be proven wrong, it is \textbf{scientific} (unlike unfalsifiable statements such as ``unemployment is a way to heaven'').
\end{enumerate}

\begin{CriticalPoint}
\textbf{Value neutrality:}
\begin{itemize}
\item \textbf{Value neutrality} = the researcher's ability to understand social action without letting values influence the research (Weber). Values may influence the \textbf{choice of topic} (value relevance), but research should be objective after selection. Later critics (Gouldner, Medel) argue objectivity is unattainable even in later stages (data interpretation, funding, etc.).
\end{itemize}
\end{CriticalPoint}

\paragraph{Q. ``Common sense is familiarization with reality; sociology is de-familiarization.''}
\begin{itemize}
\item Explain the difference between sociology and common sense (with examples). Sociology ``de-familiarises'' the everyday world; common sense simply familiarises us with it.
\end{itemize}

\begin{QuickRevision}
\begin{itemize}
\item Conservative reaction question $\rightarrow$ define modernity + revolutions + features + conservative (de Bonald/Comte/Durkheim) \& radical (Marxism) reactions.
\item Ethnomethodology $\rightarrow$ define + phenomenology + Garfinkel + breaching/documentary/indexicality.
\item Science vs common sense $\rightarrow$ sociology multi-paradigmatic; interpretative sociology bridges; Lyotard.
\item Unemployment$\rightarrow$crime $\rightarrow$ Popper (falsifiability) + positivism + non-positivist/Marxist critique $\rightarrow$ scientific.
\end{itemize}
\end{QuickRevision}

\sectiondivider

\section{Robert K. Merton --- Overview \& Deviance Typology}

\subsection{Answer-Writing: ``University as a Social System''}

The lecture opens with a PYQ: \textbf{``Critically analyze the extent to which a university is a social system.''} This is applied knowledge --- the theory (social system) is known, and it must be applied.

\begin{PYQLink}
\begin{itemize}
\item Applied knowledge is very important --- even if an applied question is not asked, one should have it. This is a previous-year question.
\end{itemize}
\end{PYQLink}

\paragraph{The approach}
\begin{enumerate}
\item \textbf{Identify keywords} --- ``social system'' is the keyword; ``university'' is the supportive keyword.
\item \textbf{Define social system}, then \textbf{describe the four needs} (functional prerequisites) = \textbf{AGIL}.
\item \textbf{Apply AGIL to the university:}
\end{enumerate}

\begin{ComparisonBox}
\textbf{University as a social system (AGIL applied):}\vspace{4pt}
\begin{center}
\rowcolors{2}{white}{Tint}
\begin{tabularx}{\textwidth}{p{2.8cm} p{3.8cm} Y}
\rowcolor{AccentDark}
\thead{Need} & \thead{Meaning} & \thead{How the university fulfils it} \\
\midrule
\textbf{Adaptation} & Surviving by adapting to the external environment & Introducing world best practices; new courses (environment, cyber security, AI); English-medium proliferation; centralised AC library; global infrastructure; online/distance education (pandemic adaptation by JNU/Delhi University) \\
\textbf{Goal attainment} & Setting and achieving goals & 100\% placement; maximum research scholars; fellowships; high NIRF/global ranking; three publications per scholar \\
\textbf{Integration} & Resolving conflict & Anti-ragging committee; proctor/rector/chancellor; student grievance-redressal cell; internal committees before cases reach the Supreme Court \\
\textbf{Latency} & Motivating, diffusing tension & Extracurricular activities, cultural events, fests (dance/music) that sustain motivation and diffuse academic tension (prevent suicide) \\
\bottomrule
\end{tabularx}
\end{center}
\end{ComparisonBox}

\begin{CriticalPoint}
\begin{itemize}
\item Always \textbf{define each term first} (adaptation, goal attainment, integration, latency), then apply it. Some items (student union, international collaboration) can be placed in more than one box, as long as it makes sense to the examiner.
\end{itemize}
\end{CriticalPoint}

\begin{QuickRevision}
\begin{itemize}
\item Social system = entity with AGIL organs; university qualifies.
\item A (new courses, online education, infrastructure), G (placement, rankings, fellowships), I (anti-ragging, grievance cells), L (fests, extracurricular).
\end{itemize}
\end{QuickRevision}

\sectiondivider

\subsection{The Context of Talcott Parsons' Writing}

Before Merton, the lecture sets the \textbf{context} of Parsons' work --- because a scholar's writing reflects his time and social location (unlike natural science, social scientists are shaped by their context).

\begin{DataFact}
\begin{itemize}
\item Around \textbf{1931}, Parsons joined \textbf{Harvard University} --- but as an \textbf{economist}, later switching to sociology.
\item His \textbf{PhD thesis was on the concept of capitalism}; he read \textbf{Sombart} (a critic of Marx on capitalism) and attended conferences on Marx conducted by \textbf{Marianne Weber}; he studied at the \textbf{London School of Economics} alongside anthropologists like \textbf{Malinowski, Morris Ginsberg and L. T. Hobhouse}.
\item The 1940s onwards was \textbf{World War II} and the start of the \textbf{Cold War} (clash of ideologies --- USA vs Soviet). America had just recovered from the \textbf{Great Depression}. India was part of the ideological clash but chose the \textbf{Non-Aligned Movement (NAM)} --- neither capitalist nor communist.
\item So a scholar like Parsons writing on \textbf{social order} was inevitable --- he was writing in, and reflecting, a time of crisis. The problem is that, while writing on American capitalism, he was \textbf{justifying the American social system as just and perfect}, with maximum adaptive capacity (defending it, since he was American).
\item America's first sociology department came in \textbf{1892} (Chicago). By the 1930s Parsons arrived as an economist, turning fully sociologist by 1940. There was a confrontation between him and \textbf{Sorokin} (a Russian-American critic of American capitalism) in the \textbf{department of social relations} (initially multi-disciplinary --- economics, anthropology, sociology, political science, philosophy --- hence ``social relations'', a model the JNU Centre for the Study of Social Systems follows).
\item Structural functionalism remained dominant \textbf{until 1970}. The \textbf{Chicago school (1892)} was empirical/pragmatic (social-reformist), while Parsons brought \textbf{grand theorisation} --- his aim was to shift the centre of sociology \textbf{from Europe to America}.
\item Structural functionalism declined by the end of the 1970s because it \textbf{could not explain conflict}. It studied equilibrium (origin and end), never conflict as such. Then \textbf{C. Wright Mills} (who introduced Marxism into America) and \textbf{Alvin Gouldner} (author of \textit{The Coming Crisis of Western Sociology}) led the critique.
\end{itemize}
\end{DataFact}

\begin{CriticalPoint}
\begin{itemize}
\item C. Wright Mills argued structural functionalism is \textbf{teleological} (explaining the \textit{function/effect} when asked the \textit{cause} --- ``a chair is to sit on'') and \textbf{tautological} (using many words to say the same thing).
\item Gouldner's ``crisis of Western sociology'' = the decline of structural functionalism with the rise of Marxism.
\end{itemize}
\end{CriticalPoint}

\begin{DataFact}
\begin{itemize}
\item \textbf{Robert Merton} was added to the UPSC syllabus in \textbf{2003} (and \textbf{Mead} was added later). Other significant living/recent sociologists: \textbf{Anthony Giddens}, \textbf{Zygmunt Bauman} (died $\sim$3 years before the lecture), \textbf{Immanuel Wallerstein}, \textbf{Jürgen Habermas}.
\end{itemize}
\end{DataFact}

\begin{CriticalPoint}
\begin{itemize}
\item University teachers are recruited through \textbf{NET} (National Eligibility Test --- a ``licence'' to teach), and \textbf{JRF} (a fellowship of \rupee 35,000/month for 5 years). The exam is easy relative to the responsibility, and this is why Indian sociology is not developing --- a book on the top 50 sociologists of the world had \textbf{no Indian} (only \textbf{Louis Dumont}, a French scholar who did fieldwork in India, featured). Marx, Weber, Durkheim, Parsons and Merton are \textbf{not field workers} --- they theorise.
\end{itemize}
\end{CriticalPoint}

\begin{QuickRevision}
\begin{itemize}
\item Parsons: economist $\rightarrow$ sociologist; PhD on capitalism; LSE (Malinowski, Ginsberg, Hobhouse); wrote in WWII/Cold War context; justified America.
\item Structural functionalism dominant till 1970; declined (couldn't explain conflict) --- critique by C. Wright Mills (teleological/tautological) and Alvin Gouldner (\textit{Coming Crisis of Western Sociology}).
\item Merton added to UPSC syllabus 2003.
\end{itemize}
\end{QuickRevision}

\sectiondivider

\subsection{Introducing Robert K. Merton}

\begin{KeyConcept}
\begin{itemize}
\item Merton was inspired by \textbf{Durkheim, Marx and Weber} and was a heavy critic of \textbf{Parsons}. So one can think of Merton as \textbf{``Marx + Weber + Durkheim $-$ Parsons.''}
\item He is \textbf{not a field worker} --- but he criticises grand theorisation, asking: how does this theory help me with everyday problems?
\end{itemize}
\end{KeyConcept}

\begin{DataFact}
\begin{itemize}
\item Merton introduced sociology at the \textbf{University of Cambridge}, and was a \textbf{student of Parsons} (he was 21 when he heard Parsons' 1931 ``Structure of Social Action'' lecture --- nobody could comprehend it).
\item Parsons was invited to Cambridge to explain sociology; people went ``mad'' (the lecture was incomprehensible). Parsons' ambition was to make sociology \textbf{as complex as physics} (so no one outside could understand it). Merton opposed this --- sociology should be \textbf{spread}, not killed. He established the Cambridge sociology department.
\end{itemize}
\end{DataFact}

\paragraph{Merton's contributions}
\begin{enumerate}
\item \textbf{Criminology} --- the study of crime (very popular in America; the world of crime fascinates because it is ``abnormal''). To understand why people become criminals he read psychology, cultural theories, Marxism and biological theories, then developed a \textbf{sociological theory of deviance} --- the \textbf{strain theory} (also called the \textbf{theory of anomie} or the \textbf{deviance typology}). The word ``anomie'' he borrowed from \textbf{Durkheim}.
\item \textbf{Sociology of science} --- he is the \textbf{founder of the sociology of science} (ethos of science, norms of science, the \textbf{CUDOS} principle), and the \textbf{first sociologist ever awarded the National Medal of Science} in America (for studying the behaviour of scientists).
\item \textbf{Bureaucratic personality} --- he studied the \textbf{dysfunction of bureaucracy} (the bureaucrat who is supposed to be goal-oriented but turns means-oriented --- ``means become ends in themselves'').
\item \textbf{Reference group behaviour} --- why we follow people on Instagram/Facebook/YouTube (Marxism, Durkheim, Weber, Parsons cannot explain this everyday reality).
\end{enumerate}

\begin{CriticalPoint}
\begin{itemize}
\item Merton's interest grew from the \textbf{Chicago school} (empiricism, pragmatism, everyday reality). The Chicago school studied \textbf{social issues} (crime rate, poverty, racism, homelessness).
\item \textbf{Public sociology} (e.g., \textbf{Jane Addams}, social activism/charity, helping the poor --- later developed ``public sociology'') is oriented to \textit{developing society}, not theories. \textbf{Professional sociology} (Marx, Weber, Durkheim) develops \textit{theories}. This distinction is by \textbf{Michael Burawoy}.
\item \textbf{W. E. B. Du Bois} (dual consciousness) studied racism.
\end{itemize}
\end{CriticalPoint}

\begin{QuickRevision}
\begin{itemize}
\item Merton = Marx + Weber + Durkheim $-$ Parsons; critic of grand theory; not a field worker.
\item Contributions: criminology (strain theory/anomie), sociology of science (founder; CUDOS; first sociologist with National Medal of Science), bureaucratic personality (dysfunction), reference group behaviour.
\item Public vs professional sociology (Burawoy); Chicago school = empiricism.
\end{itemize}
\end{QuickRevision}

\sectiondivider

\subsection{Merton's Predecessors: The Chicago School}

Before Merton's own theory, three Chicago-school scholars studied crime, and Merton is related to them:

\begin{DataFact}
\begin{itemize}
\item Other Chicago scholars: \textbf{Jane Addams} (public sociology), \textbf{Du Bois} (dual consciousness), \textbf{Small, Lester Ward, Sumner}, and later \textbf{W. I. Thomas}, \textbf{Charles Horton Cooley}, and \textbf{George Herbert Mead} (in this syllabus).
\end{itemize}
\end{DataFact}

\paragraph{W. I. Thomas --- the Thomas theorem}

\begin{KeyConcept}
\begin{itemize}
\item \textbf{``When you define something as real, it is real in its consequences.''} If you tell a student ``you cannot crack UPSC,'' he will believe it, lose motivation, and eventually leave preparation; if you say ``you will make it,'' he stays motivated all day.
\item Applied to crime: when you \textbf{label} people as criminals, they \textbf{become} criminals. (Black people were labelled as criminals, and their self-perception became criminal.)
\end{itemize}
\end{KeyConcept}

\paragraph{Charles Horton Cooley --- the looking-glass self}

\begin{KeyConcept}
\begin{itemize}
\item \textbf{``I am not what I think I am; I am not what you think I am; I am what I think you think I am.''}
\item You behave towards someone (e.g., an ex-girlfriend, or a teacher) according to what you think \textit{they} think of you. If 10 people call you cool, you start believing you are cool --- your ``coolness'' is the result of your \textit{perception} that they think you are cool.
\item \textbf{Looking-glass self:} after an interaction, you enact it in front of a mirror, imagining what the other person thought of you --- your self-image is a reflection (mirror) of that perception.
\end{itemize}
\end{KeyConcept}

\paragraph{George Herbert Mead}
Mead is \textbf{much influenced by Charles Horton Cooley} (detailed later in the syllabus).

\begin{KeyConcept}
\textbf{Self-fulfilling prophecy:}
\begin{itemize}
\item Inspired by Thomas but \textbf{coined by Merton}. The same idea applied to large-scale social phenomena: if a rumour spreads that \textbf{ICICI Bank} is going to collapse, everyone rushes to withdraw, and the bank actually collapses --- \textbf{Yes Bank}'s collapse was also a self-fulfilling prophecy (not purely economic).
\item Merton applied this to education: \textbf{inequality in education / dropouts} can be explained by self-fulfilling prophecy (a teacher who demoralises students causes dropouts, which economists would wrongly explain only by distance/infrastructure).
\item Note: the self-fulfilling/labelling approach explains \textit{black} crime, but \textbf{not white crime} --- which is why Merton moved to his own theory.
\end{itemize}
\end{KeyConcept}

\begin{QuickRevision}
\begin{itemize}
\item \textbf{Thomas theorem:} ``if you define something as real, it is real in its consequences.''
\item \textbf{Cooley's looking-glass self:} ``I am what I think you think I am.''
\item \textbf{Self-fulfilling prophecy} (coined by Merton, inspired by Thomas) --- ICICI/Yes Bank bank-run example.
\item Mead influenced by Cooley.
\end{itemize}
\end{QuickRevision}

\sectiondivider

\subsection{Theories of Deviance Reviewed by Merton}

Merton reviewed several existing theories of deviance (his \textbf{predecessors}) and found each incomplete for his question --- \textit{why is the crime rate rising?}

\begin{enumerate}
\item \textbf{Psychological theory} --- Cooley, Thomas, and \textbf{Freud} (id/ego/superego; criminals have an over-developed id --- impulsive, no balance).
\item \textbf{Biological theory} --- \textbf{Lombroso} (\textit{Born Criminals}): criminals have unique physiological features (long nose, broad forehead, elongated ears). (Merton's point: physiology differs, but that does not explain \textit{why} someone becomes criminal.)
\item \textbf{Marxist theory} --- \textbf{capitalism is criminogenic} (its values reproduce criminal behaviour: appropriation of surplus value/wealth; billionaires are celebrated). But it cannot explain \textbf{why the working class} (not the bourgeoisie) commit street crime and go to jail.
\item \textbf{Interactionist theory} --- \textbf{Howard Becker's labelling theory} (you define someone as criminal, they become criminal; but crime is contextual --- the same act may not be crime elsewhere).
\item \textbf{Differential association theory} --- \textbf{Sutherland}: you become criminal if you associate with criminals (differential association).
\item \textbf{Structural functional theory} --- deviance creates solidarity (Nirbhaya united India; 9/11 united the world in ``war on terror'') and brings social change --- so deviance is functional.
\end{enumerate}

\begin{ExampleBox}
\begin{itemize}
\item If you murder someone wearing an army uniform (and the victim is an enemy), you are a hero, not a criminal. The same act in another context is crime --- the uniform (a symbol) legitimises the killing. This shows ``legitimate monopoly of violence'' and that crime is context-dependent.
\end{itemize}
\end{ExampleBox}

\begin{DataFact}
\begin{itemize}
\item Howard Becker's \textbf{labelling theory} = interactionist theory of deviance. It is close to micro-sociology (understanding from the actor's perspective in a specific historical/social/cultural context).
\end{itemize}
\end{DataFact}

\begin{QuickRevision}
\begin{itemize}
\item Six reviewed theories: psychological (Freud id/ego/superego), biological (Lombroso), Marxist (capitalism criminogenic), interactionist (Becker labelling), differential association (Sutherland), structural functional (deviance $\rightarrow$ solidarity/change).
\item None answered ``why is crime rate rising in 1940s America?'' $\rightarrow$ Merton's own theory.
\end{itemize}
\end{QuickRevision}

\sectiondivider

\subsection{Merton's Strain Theory / Deviance Typology}

\begin{KeyConcept}
\begin{itemize}
\item \textbf{Anomie is rooted in social structure} --- it is the product of the \textbf{disjunction (mismatch) between culturally defined goals and socially legitimate means}.
\item In society there are two things: \textbf{culturally defined goals} (CDG) and \textbf{socially legitimate means} (SLM). When there is a \textbf{mismatch} between them, \textbf{strain} develops.
\end{itemize}
\end{KeyConcept}

\begin{KeyConcept}
\begin{itemize}
\item The culturally defined goal (earn money, gain status, own a car/bungalow, fame, wealth) Merton calls the \textbf{American dream} --- it came from America via \textbf{Hollywood} (whose movies glorified the individual who struggled and earned billions; heroes = white/money-makers, villains = black/money-stealers). Globalisation is largely Americanisation; our rising aspirations follow the American trend.
\end{itemize}
\end{KeyConcept}

\begin{DataFact}
\begin{itemize}
\item To understand deviance, we must understand the \textbf{structural conditions} of society --- i.e., the \textbf{generalised values} internalised by actors (that ``this is what success is'').
\item Example: \textbf{Al Capone} (the ``Dawood Ibrahim'' of America) --- wealth, car, bungalow without education/hard work, via illegitimate means, yet on the NYPD's ``most wanted'' list.
\end{itemize}
\end{DataFact}

\paragraph{The five types of adaptation}

\begin{center}
\rowcolors{2}{white}{Tint}
\begin{tabularx}{\textwidth}{p{2.6cm} p{2.6cm} p{2.6cm} p{1.8cm} Y}
\rowcolor{AccentDark}
\thead{Type} & \thead{Goals (CDG)} & \thead{Means (SLM)} & \thead{Deviant?} & \thead{Example} \\
\midrule
\textbf{Conformist} & + (accept) & + (accept) & No & Affluent/privileged (can afford coaching, everything) \\
\textbf{Innovator} & + (accept) & $-$ (reject) & Yes & Al Capone, drug peddlers, money launderers, Pablo Escobar \\
\textbf{Ritualist} & $-$ (reject) & + (accept) & Yes & Indian caste structure; low-grade bureaucrats (means become ends) \\
\textbf{Retreatist} & $-$ (reject) & $-$ (reject) & No & Hippies, vagabonds, ascetics \\
\textbf{Rebellion} & $-$/+ (new goals) & $-$/+ (new means) & Yes & Gandhi, Hitler, Marx, Ambedkar \\
\bottomrule
\end{tabularx}
\end{center}

\begin{KeyConcept}
\begin{itemize}
\item Deviance = mismatch between CDG and SLM. So \textbf{only two groups are not deviant} (conformist +/+, retreatist $-/-$); \textbf{three are deviant} (innovator, ritualist, rebellion).
\end{itemize}
\end{KeyConcept}

\begin{ExampleBox}
\textbf{Examples of each type:}
\begin{itemize}
\item \textbf{Conformist:} born into an educated middle/upper-middle-class family $\rightarrow$ socialisation $\rightarrow$ good qualification $\rightarrow$ good placement.
\item \textbf{Innovator:} subscribes to goals (car, bungalow, luxury) but uses illegitimate means (selling illegal substances; money laundering).
\item \textbf{Ritualist:} not concerned with material goals; follows legitimate means for their own sake. The \textbf{Indian caste structure} is based on ritualism (lower castes' dharma is menial work --- no aspirations). \textbf{Low-grade bureaucrats} are ritualist (means-oriented, not goal-oriented; ``means become an end in itself'' --- the dysfunction of bureaucracy).
\item \textbf{Retreatist:} neither goals nor means --- hippies, vagabonds (always on the beach, no goal).
\item \textbf{Rebellion:} rejects CDG and SLM, but substitutes \textbf{new goals and new means} they want to make society's new norms. Inspired by Weber's \textbf{charismatic authority} (extraordinary qualities, bring change). Gandhi, Hitler, Marx, Ambedkar --- all rebellions (Gandhi's ``Ram Rajya'': limited needs, trusteeship, village panchayat/commune system; Marx's revolution).
\end{itemize}
\end{ExampleBox}

\begin{CriticalPoint}
\begin{itemize}
\item ``Rebellion'' is a socially constructed label --- one person's rebel is another's freedom fighter (Gandhi was a rebel to the British, a leader to his followers).
\end{itemize}
\end{CriticalPoint}

\begin{DataFact}
\begin{itemize}
\item \textbf{Latif Chaudhary} studied South Asian bureaucracies (Nepal, Bhutan, Bangladesh, India, Sri Lanka --- highly corrupt in the corruption index) and labelled them \textbf{rent-seeking bureaucracies} (rent-seekers, not efficient). Such bureaucrats can fit both the \textbf{ritualist} and \textbf{innovator} categories.
\end{itemize}
\end{DataFact}

\begin{KeyConcept}
\begin{itemize}
\item The culturally defined goal (money/status) is available to everyone (democracy promises equal opportunity), but the \textbf{legitimate means are not equally available} (because of social location --- poverty, caste/class discrimination, transgender exclusion). Everyone lives under the pressure to have money/fame, so those who cannot get it legitimately resort to \textbf{illegitimate means}. This is what explains crime --- not poverty alone.
\item Note: the \textbf{Ministry of Social Justice} wrote to UPSC that there should be no caste-based discrimination in the personality test (caste column stripped off) --- evidence of caste-based barriers.
\end{itemize}
\end{KeyConcept}

\begin{QuickRevision}
\begin{itemize}
\item Anomie rooted in \textbf{social structure} = disjunction between \textbf{CDG} and \textbf{SLM}.
\item \textbf{American dream} (from Hollywood) = CDG.
\item Five types: conformist (+/+), innovator (+/$-$), ritualist ($-$/+), retreatist ($-/-$), rebellion ($-$/+ new).
\item Only conformist \& retreatist are non-deviant; innovator, ritualist, rebellion are deviant.
\item Crime rises because goals are universal but means are unequally distributed by social location.
\end{itemize}
\end{QuickRevision}

\sectiondivider

\section{Robert K. Merton --- Reference Group, Matthew Effect \& Middle-Range Theory}

\subsection{Answer-Writing: ``AGIL \& Key Problems in Society'' (UPSC 2018)}

The lecture opens with a 20-mark UPSC question (2018): \textbf{``How can Parsons' AGIL framework be used to analyze key problems in a society? Discuss.''}

\begin{PYQLink}
\begin{itemize}
\item This is a \textbf{20-mark UPSC (2018)} question. Time budget: $\sim$10--12 minutes; structure matters.
\end{itemize}
\end{PYQLink}

\paragraph{The approach}
\begin{enumerate}
\item \textbf{Define the AGIL framework} --- define A, G, I, L; explain they are the \textit{needs of society}; if needs are fulfilled, society survives (equilibrium). Give examples and systems (adaptation $\rightarrow$ economy/organismic system, etc.).
\item \textbf{State the paradigm} --- this is the \textbf{structural-functional paradigm} developed in America by Talcott Parsons in his book \textit{The Social System} (1951).
\item \textbf{Introduce ``key problems''} --- poverty, unemployment, crime against women, casteism, migration (especially forced/push-factor migration).
\item \textbf{Show each key problem = failure of a system:}
\begin{itemize}
\item \textbf{Poverty / unemployment} $\rightarrow$ failure of the \textbf{organismic system (economy)} --- production/distribution not happening; unemployment is now \textbf{pathological} (not normal).
\item \textbf{Crime against women} (non-recognition of marital rape, domestic violence) $\rightarrow$ failure of the \textbf{personality, social and cultural systems} (family, laws, judiciary).
\item \textbf{Casteism} $\rightarrow$ operates in \textbf{implicit/new form} (new casteism); shows Indians have not internalised pattern variable B (achievement).
\item \textbf{Migration} $\rightarrow$ failure of the organismic system (no local employment/resources).
\item \textbf{COVID/health} $\rightarrow$ failure of the personality (political) and organismic systems (no vaccination, awareness, infrastructure).
\item \textbf{Religion} $\rightarrow$ if religion is so functional (value consensus, motivation, tonic of self-confidence in crisis), how explain hate speech, communal violence, fundamentalism?
\end{itemize}
\end{enumerate}

\begin{CriticalPoint}
\begin{itemize}
\item A Marxist scholar asks this question. If you simply re-explain AGIL, you have \textbf{not answered} --- you must use AGIL to \textit{analyse} the problems (show how each problem is a failure of a system).
\item \textbf{Always discuss government steps/solutions} too (e.g., NEP 2020 for education) --- take the side of government while highlighting problems.
\end{itemize}
\end{CriticalPoint}

\begin{DataFact}
\begin{itemize}
\item Parsons remarked ``who now reads Spencer?'' (actually said by \textbf{Brinton}, but Parsons agreed). Spencer became irrelevant because, after the 1930s, the state had to intervene (no more laissez-faire).
\item \textbf{C. Wright Mills} replied, ``who now reads Parsons?'' --- marking the crisis of Western sociology and the shift to Marxism (Frankfurt School, neo-Marxists).
\end{itemize}
\end{DataFact}

\begin{QuickRevision}
\begin{itemize}
\item Define AGIL $\rightarrow$ state structural-functionalism (\textit{Social System}, 1951) $\rightarrow$ key problems $\rightarrow$ map each to a system's failure $\rightarrow$ add government measures.
\item ``Who now reads Spencer?'' (Brinton/Parsons) $\leftrightarrow$ C. Wright Mills: ``who now reads Parsons?''
\end{itemize}
\end{QuickRevision}

\sectiondivider

\subsection{Merton's Deviance Typology: The Five Types}

\begin{DataFact}
\textbf{The six reviewed theories (revision):}
\begin{itemize}
\item Merton reviewed six theories before his own: (1) psychological (Freud id/ego/superego), (2) biological (Lombroso), (3) Marxist (capitalism criminogenic), (4) interactionist (Howard Becker labelling), (5) differential association (Sutherland), (6) structural functional (deviance $\rightarrow$ solidarity/change).
\end{itemize}
\end{DataFact}

\begin{KeyConcept}
\begin{itemize}
\item \textbf{Anomie is rooted in the social structure} and is the product of the \textbf{disjunction between culturally defined goals (CDG) and socially legitimate means (SLM)}.
\item Deviance arises from \textbf{mismatch}. So only the conformist (+/+) and retreatist ($-/-$) are \textbf{not} deviant; the innovator (+/$-$), ritualist ($-$/+) and rebellion ($-$/+) are \textbf{deviant}.
\end{itemize}
\end{KeyConcept}

\begin{ExampleBox}
\begin{itemize}
\item \textbf{Conformist:} born to an affluent family, socialised to achieve (coaching, good colleges, placements).
\item \textbf{Innovator:} accepts the goal (wealth) but uses illegitimate means --- drug peddlers, money launderers.
\item \textbf{Ritualist:} Indian caste structure (lower castes have no material aspirations --- their dharma is menial work); low-grade bureaucrats (means-oriented, red-tapism).
\item \textbf{Retreatist:} hippies, vagabonds.
\item \textbf{Rebellion:} Gandhi (Ram Rajya, trusteeship, village panchayats/commune, stateless society), Marx (revolution), Ambedkar, Hitler, Nelson Mandela, Bhagat Singh --- all ``charismatic'' figures with a new vision.
\item \textbf{Escorts (highly educated)} are \textbf{innovators} --- they accept the goal but lack legitimate means.
\end{itemize}
\end{ExampleBox}

\begin{CriticalPoint}
\begin{itemize}
\item Nobody is born a conformist, innovator, ritualist, retreatist or rebellion (Gandhi was not born a rebel; hippies not born retreatists). The \textbf{structures of society shape them} --- it is not personal trouble but a \textbf{public issue} (if it were personal, it would be limited to one or two cases; seen on a large scale, it is public).
\end{itemize}
\end{CriticalPoint}

\begin{QuickRevision}
\begin{itemize}
\item Six reviewed theories $\rightarrow$ Merton's own: anomie rooted in social structure; disjunction CDG vs SLM.
\item Non-deviant: conformist (+/+), retreatist ($-/-$). Deviant: innovator (+/$-$), ritualist ($-$/+), rebellion (new goals/means).
\item Types are produced by social structure, not born (personal trouble vs public issue).
\end{itemize}
\end{QuickRevision}

\sectiondivider

\subsection{The Matthew Effect \& Relative Deprivation}

\begin{KeyConcept}
\textbf{Matthew effect:}
\begin{itemize}
\item Very few people become criminals; most are conformists --- because conformists have a \textbf{comparative advantage}. Merton's point: \textbf{``advantage begets advantage''} (the Matthew effect = comparative advantage).
\item Marx and Benjamin Franklin said ``money begets money'' (M1 $\rightarrow$ M2). Merton says it is \textbf{not money that begets money, but advantage that begets advantage} --- being placed advantageously brings more advantage.
\end{itemize}
\end{KeyConcept}

\begin{ExampleBox}
\begin{itemize}
\item \textbf{DNA discovery:} \textbf{Rosalind Franklin} was the actual contributor, but the (male) scientists got the Nobel Prize --- being male is a comparative advantage.
\item \textbf{Abhijit Banerjee} became more famous than his wife (also an economist) --- a male comparative advantage (invited to more conferences).
\item \textbf{Rahul Gandhi} studied at Oxford/Harvard not merely because of money/qualification (there are richer people in India) but because he belongs to the \textbf{Gandhi family} --- advantage begets advantage.
\end{itemize}
\end{ExampleBox}

\begin{KeyConcept}
\textbf{Relative deprivation (vs absolute deprivation):}
\begin{itemize}
\item When you meet someone in an advantageous position, you feel \textbf{inferior/deprived} --- this is \textbf{relative deprivation} (relative alienation).
\item Marx talks about \textbf{haves and have-nots} (absolute deprivation --- owners vs non-owners of means of production). Merton says there is \textbf{no absolute deprivation, only relative} --- even among the proletariat, some are richer than others; you compare yourself to your \textit{reference point}, not to Ambani or Tata.
\end{itemize}
\end{KeyConcept}

\begin{ExampleBox}
\textbf{COVID \& social distancing:}
\begin{itemize}
\item \textbf{Social distancing is a privilege} --- only the advantageously placed (big houses, running water, sanitation) can maintain it; the homeless/one-bedroom families cannot. The people who contracted corona were also \textbf{relatively deprived} people (those who got two vaccine doses early had a comparative advantage --- the ``Matthew effect'').
\end{itemize}
\end{ExampleBox}

\begin{QuickRevision}
\begin{itemize}
\item \textbf{Matthew effect} = comparative advantage; ``advantage begets advantage'' (vs Marx/Franklin's ``money begets money'').
\item Examples: Rosalind Franklin (DNA), Abhijit Banerjee, Rahul Gandhi (family advantage).
\item \textbf{Relative deprivation} (not absolute --- vs Marx's haves/have-nots); social distancing = privilege.
\end{itemize}
\end{QuickRevision}

\sectiondivider

\subsection{Reference Group Theory}

\begin{KeyConcept}
\begin{itemize}
\item Every person compares/evaluates himself in relation to people he thinks are well-off or well-placed. \textbf{Reference group theory} explains how individuals/groups \textbf{evaluate themselves in relation to other individuals or groups}.
\item Merton also used the term \textbf{reference individual} --- later called a \textbf{role model} (the kind of person I want to be like).
\end{itemize}
\end{KeyConcept}

\begin{DataFact}
\begin{itemize}
\item \textbf{Membership group} (also called \textbf{in-group}, a term from psychologist \textbf{Sumner}) --- the group you belong to.
\item \textbf{Non-membership group} (\textbf{out-group}) --- the group you want to join.
\end{itemize}
\end{DataFact}

\begin{ExampleBox}
\textbf{The cricketer:}
\begin{itemize}
\item If I am a cricketer in the Ranji Trophy team, my Ranji team is my \textbf{membership group}; the T20 World Cup team is my \textbf{non-membership group} --- I want to be there, and I feel \textbf{relative deprivation}.
\item I resolve this by \textbf{anticipatory socialisation} --- doing everything I think will get me membership (talking to team leaders, building a routine) until I finally enter the T20 team.
\end{itemize}
\end{ExampleBox}

\begin{KeyConcept}
\begin{itemize}
\item \textbf{Anticipatory socialisation is functional in open societies} (meritocracies --- it enables social mobility) and \textbf{dysfunctional in closed societies} (no equal opportunity --- climbing the ladder brings consequences).
\end{itemize}
\end{KeyConcept}

\begin{ExampleBox}
\textbf{Open vs closed society:}
\begin{itemize}
\item \textbf{Open society (functional):} a bright student gets into a prestigious college on merit $\rightarrow$ anticipatory socialisation $\rightarrow$ social mobility.
\item \textbf{Closed society (dysfunctional):} a Shudra emulating a Brahmin (wearing the sacred thread, the tilak, reading the Vedas) is rejected by both Brahmin and Shudra communities --- he becomes a \textbf{marginal man} (a term from \textbf{Robert Ezra Park}, Chicago school). The \textbf{marginal man is the product of the failure of anticipatory socialisation}.
\item \textbf{Burgess} coined the term \textbf{social distance} (and the \textbf{social distance scale}, developed to measure white attitudes towards blacks).
\end{itemize}
\end{ExampleBox}

\begin{CriticalPoint}
\textbf{Vertical vs horizontal untouchability:}
\begin{itemize}
\item Untouchability was once \textbf{vertical} (Brahmin vs Dalit/Shudra --- and it was mutual: Brahmins were impure to Shudras too; labelling only Brahmins ``pure'' means internalising Brahminical ideology).
\item Today it is \textbf{horizontal} --- from a \textbf{caste body} we have become a \textbf{bio-body} (concerned with health/COVID purity, not caste pollution). From \textbf{vertical untouchability} we have moved to \textbf{horizontal untouchability}.
\end{itemize}
\end{CriticalPoint}

\paragraph{What reference group theory explains}

\begin{DataFact}
\begin{itemize}
\item \textbf{Social mobility} --- the auto-rickshaw driver's son who becomes an IAS officer (he felt relative deprivation, underwent anticipatory socialisation, and jumped vertically).
\item \textbf{Mass communication, PR, marketing of political leaders} --- and the \textbf{fashion industry} (you see a fancy pair of jeans, find out the brand, and buy the same one, even second-hand from eBay).
\item \textbf{International relations} --- the third world feels relative deprivation vis-à-vis the first world, hence demands technology transfer.
\item \textbf{Women's labour force participation} --- Indian women's participation is falling compared to global levels, so women feel relatively deprived $\rightarrow$ pressure groups demand employment.
\item \textbf{Wages for domestic work} --- domestic workers are relatively deprived vis-à-vis male breadwinners $\rightarrow$ vulnerable to abuse.
\item \textbf{Maratha reservation demand} --- the Marathas (a \textbf{dominant caste}) feel relatively deprived; dominant castes include Yadav (UP), Patidar/Patel (Gujarat), Jat (Haryana), Ahir (western UP).
\item \textbf{Kshatriya-isation} --- Shudras who claimed Kshatriya status (via Jati Purans, anticipatory socialisation) to gain higher status; the \textbf{Yadav vote bank} is also a case.
\end{itemize}
\end{DataFact}

\begin{CriticalPoint}
\begin{itemize}
\item Democracy is not ``majority'' --- it is the \textbf{``tyranny of the majority.''} SCs/STs are few in number, so the majority (upper castes) can dominate. Dominant castes rule pockets across India.
\end{itemize}
\end{CriticalPoint}

\begin{QuickRevision}
\begin{itemize}
\item Reference group = the group/individual you compare yourself to; \textbf{reference individual = role model}.
\item \textbf{Membership (in-group)} vs \textbf{non-membership (out-group)}; \textbf{anticipatory socialisation} $\rightarrow$ social mobility.
\item Anticipatory socialisation functional in \textbf{open}, dysfunctional in \textbf{closed} societies $\rightarrow$ \textbf{marginal man} (Robert Park).
\item Social distance (Burgess); vertical $\rightarrow$ horizontal untouchability; explains fashion, PR, IR, reservations, social mobility.
\end{itemize}
\end{QuickRevision}

\sectiondivider

\subsection{Forms of Theories: Case, Middle-Range \& Grand}

\begin{center}
\rowcolors{2}{white}{Tint}
\begin{tabularx}{\textwidth}{p{2.6cm} p{2.2cm} p{3.6cm} p{2.4cm} Y}
\rowcolor{AccentDark}
\thead{Type} & \thead{Range} & \thead{Explains} & \thead{Testability} & \thead{Examples} \\
\midrule
\textbf{Case theory} & Low range & Individual behaviour (why \textit{he} fell asleep, stole, lied) & Testable (low/zero abstraction) & Case studies of a single person \\
\textbf{Middle-range theory} & Middle range & \textbf{Group behaviour} (why women drop out more than men) & Testable & Reference group, relative deprivation, Matthew effect \\
\textbf{Grand theory} & High range & \textbf{Structures and changes in structures} & \textbf{Non-testable} (high abstraction) & Structural functionalism, Marxism, feminism, Freud's psychoanalysis, Weber's rationalisation \\
\bottomrule
\end{tabularx}
\end{center}

\begin{KeyConcept}
\textbf{What each level answers:}
\begin{itemize}
\item \textbf{Case theory:} why did he fall asleep? why did they break up? (individual)
\item \textbf{Grand theory:} why does economy support polity? how is religion functional to the system? how does the superstructure justify the economic base? (structures)
\item \textbf{Middle-range theory:} why do certain \textit{groups} behave this way? (groups) --- e.g., why women drop out of school more than men.
\end{itemize}
\end{KeyConcept}

\begin{ExampleBox}
\begin{itemize}
\item When Durkheim compared \textbf{Catholics vs Protestants} and \textbf{men vs women} (suicide), he was at the \textbf{middle range}. When he generalised ``with industrialisation, suicide rates increase,'' he moved to the \textbf{grand level}.
\end{itemize}
\end{ExampleBox}

\begin{KeyConcept}
\textbf{Grand theories and their bases:}
\begin{itemize}
\item \textbf{Marx} $\rightarrow$ base is economy; \textbf{Freud} $\rightarrow$ base is sex (society as the result of sex prohibition/taboo); \textbf{Weber} $\rightarrow$ rationalisation.
\item Grand theories are \textbf{highly abstract and non-testable} (no instrument to ``see'' revolution or measure that religion is 90\% functional). Middle-range theories are \textbf{testable} (e.g., the unemployment--crime-rate question from Test 1 was a middle-range theory, and it was shown to be testable).
\end{itemize}
\end{KeyConcept}

\begin{CriticalPoint}
\textbf{Merton's programme:}
\begin{itemize}
\item As an \textbf{empirical sociologist}, one should have restraint: develop \textbf{middle-range concepts} to explain group behaviour, not grand narratives. Merton's own theories (reference group, relative deprivation, Matthew effect) are \textbf{middle-range theories} --- applicable and testable across time and space.
\item (Left for the next class: \textbf{manifest and latent functions}, and the \textbf{critique to Merton} --- e.g., that culturally defined goals are \textit{also} not equally distributed, and that ``becoming a criminal'' is not easy --- it requires the \textbf{illegitimate opportunity structure}, a concept from \textbf{Cloward and Ohlin}.)
\end{itemize}
\end{CriticalPoint}

\begin{DataFact}
\begin{itemize}
\item \textbf{Samir Amin} (an Egyptian Marxist, who died in 2017) gave the theory of \textbf{unequal (uneven) development} --- a notable Egyptian name in sociology.
\end{itemize}
\end{DataFact}

\begin{QuickRevision}
\begin{itemize}
\item Three theory types: \textbf{case} (low, individual, testable), \textbf{middle-range} (group, testable), \textbf{grand} (structures, non-testable).
\item Grand: structural functionalism, Marxism, feminism, Freud, Weber's rationalisation. Middle-range: reference group, relative deprivation, Matthew effect.
\item Merton's own theories are middle-range (testable, applicable across time/space).
\end{itemize}
\end{QuickRevision}